
\documentclass{article}
\usepackage{colortbl}
\usepackage{makecell}
\usepackage{multirow}
\usepackage{supertabular}

\begin{document}

\newcounter{utterance}

\centering \large Interaction Transcript for game `clean\_up', experiment `1\_medium\_7obj\_pl', episode 0 with Apertus{-}v1.5{-}8B.
\vspace{24pt}

{ \footnotesize  \setcounter{utterance}{1}
\setlength{\tabcolsep}{0pt}
\begin{supertabular}{c@{$\;$}|p{.15\linewidth}@{}p{.15\linewidth}p{.15\linewidth}p{.15\linewidth}p{.15\linewidth}p{.15\linewidth}}
    \# & $\;$A & \multicolumn{4}{c}{Game Master} & $\;\:$B\\
    \hline

    \theutterance \stepcounter{utterance}  
    & & \multicolumn{4}{p{0.6\linewidth}}{
        \cellcolor[rgb]{0.9,0.9,0.9}{
            \makecell[{{p{\linewidth}}}]{
                \texttt{\tiny{[P1$\langle$GM]}}
                \texttt{Jestem mistrzem gry, a Państwo grają w grę kooperacyjną z następującą siatką jako planszą:} \\
\\ 
\texttt{```} \\
\texttt{ 1234567} \\
\texttt{╔════╤══╗} \\
\texttt{║◌◌◌O│PM║ 1} \\
\texttt{║◌◌T◌├──╢ 2} \\
\texttt{║◌◌◌◌│◌◌║ 3} \\
\texttt{║S┌──┼──╢ 4} \\
\texttt{║◌│◌A│◌◌║ 5} \\
\texttt{╟─┴┬─┤◌◌║ 6} \\
\texttt{║◌◌│◌│D◌║ 7} \\
\texttt{╚══╧═╧══╝} \\
\\ 
\texttt{```} \\
\\ 
\texttt{* Górna krawędź wyświetla współrzędne x rosnące w prawo, a prawa krawędź współrzędne y rosnące w dół.} \\
\texttt{* Następujące obiekty są losowo umieszczone na siatce: 'P', 'O', 'T', 'S', 'D', 'A', 'M'.} \\
\\ 
\\ 
\texttt{Drugi gracz widzi wariant planszy, na której obiekty są umieszczone w różnych losowych miejscach. Nie można zobaczyć planszy drugiego gracza, a on nie może zobaczyć Państwa planszy.} \\
\\ 
\\ 
\texttt{**Cel:**} \\
\\ 
\texttt{Obaj gracze muszą przesuwać obiekty na swoim tle tak, aby identyczne obiekty znalazły się na tych samych współrzędnych. Muszą Państwo komunikować się z drugim graczem, aby uzgodnić wspólną konfigurację docelową.} \\
\\ 
\\ 
\texttt{**Zasady:**} \\
\\ 
\texttt{* W każdej turze można wysłać dokładnie jedno z dwóch poniższych poleceń:} \\
\texttt{ 1. `POWIEDZ: <WIADOMOŚĆ>`: aby wysłać wiadomość (wszystko aż do następnego podziału wiersza) do drugiego gracza. Przekażę ją Państwa partnerowi.} \\
\texttt{2. `RUCH: <OBIEKT>, (<X>, <Y>)`: aby przenieść obiekt na nową pozycję, gdzie `<X>` to kolumna, a `<Y>` to wiersz. Poinformuję Państwa, czy przeniesienie się powiodło, czy nie.} \\
\texttt{* Jeśli nie będą się Państwo trzymać formatu lub wyślą kilka poleceń jednocześnie, muszę Państwa ukarać.} \\
\texttt{* Jeśli obaj gracze zgromadzą więcej niż 16 kar, obaj przegrają grę.} \\
\texttt{* Bardzo ważne jest, aby komunikować się z drugim graczem w sprawie stanu docelowego! *Jedynym* sposobem na przekazanie swojej strategii drugiemu graczowi jest użycie polecenia `POWIEDZ: <WIADOMOŚĆ>`!} \\
\\ 
\\ 
\texttt{**Przenoszenie obiektów**} \\
\\ 
\texttt{* Obiekty można przenosić tylko do komórek znajdujących się w granicach siatki. Komórka docelowa musi być pusta, tzn. musi zawierać tylko symbol '◌'.} \\
\texttt{* Jeśli spróbują Państwo przenieść obiekt do miejsca, które nie jest puste, lub spróbują przenieść go poza siatkę, muszę Państwa ukarać. Otrzymają Państwo kolejną próbę.} \\
\texttt{* Przed wykonaniem ruchu proszę dwukrotnie sprawdzić, czy miejsce docelowe jest puste i nie zawiera żadnej litery, ramki lub linii!} \\
\\ 
\\ 
\texttt{**Koniec gry**} \\
\\ 
\texttt{Jeśli uważają Państwo, że osiągnęli cel wyrównania wszystkich obiektów, mogą Państwo poprosić drugiego gracza o zakończenie gry, wysyłając `POWIEDZ: zakończono?`. Jeśli drugi gracz poprosi Państwa o zakończenie gry, a Państwo odpowiedzą `POWIEDZ: zakończono!`, gra się zakończy.} \\
\\ 
\texttt{Obaj gracze wygrywają, jeśli gra zakończy się w ciągu 28 rund, gdzie jedna runda jest definiowana jako wysłanie ważnego polecenia przez dwóch graczy.} \\
\\ 
\texttt{**Punktacja:**} \\
\\ 
\texttt{Im bliżej siebie znajdują się identyczne obiekty na obu planszach, tym więcej punktów otrzymują Państwo. Kary zmniejszają liczbę punktów. Czy uda się Państwu pobić rekord?} \\
\\ 
\\ 
\texttt{Proszę wysłać wiadomość do drugiego gracza, aby rozpocząć grę!} \\
            }
        }
    }
    & \\ \\

    \theutterance \stepcounter{utterance}  
    & \multicolumn{4}{p{0.6\linewidth}}{
        \cellcolor[rgb]{0.9,0.9,0.9}{
            \makecell[{{p{\linewidth}}}]{
                \texttt{\tiny{[P1$\rangle$GM]}}
                \texttt{POWIEDZ: Gotowy do rozpoczęcia gry!} \\
            }
        }
    }
    & & \\ \\

    \theutterance \stepcounter{utterance}  
    & & & \multicolumn{2}{p{0.3\linewidth}}{
        \cellcolor[rgb]{0.9,0.9,0.9}{
            \makecell[{{p{\linewidth}}}]{
                \texttt{\tiny{[GM$|$GM]}}
                \texttt{POWIEDZ: Gotowy do rozpoczęcia gry!} \\
            }
        }
    }
    & & \\ \\

    \theutterance \stepcounter{utterance}  
    & & \multicolumn{4}{p{0.6\linewidth}}{
        \cellcolor[rgb]{0.9,0.9,0.9}{
            \makecell[{{p{\linewidth}}}]{
                \texttt{\tiny{[GM$\rangle$P2]}}
                \texttt{Jestem mistrzem gry, a Państwo grają w grę kooperacyjną z następującą siatką jako planszą:} \\
\\ 
\texttt{```} \\
\texttt{ 1234567} \\
\texttt{╔════╤══╗} \\
\texttt{║◌◌◌◌│O◌║ 1} \\
\texttt{║P◌M◌├──╢ 2} \\
\texttt{║◌S◌◌│T◌║ 3} \\
\texttt{║◌┌──┼──╢ 4} \\
\texttt{║A│◌◌│◌◌║ 5} \\
\texttt{╟─┴┬─┤◌D║ 6} \\
\texttt{║◌◌│◌│◌◌║ 7} \\
\texttt{╚══╧═╧══╝} \\
\\ 
\texttt{```} \\
\\ 
\texttt{* Górna krawędź wyświetla współrzędne x rosnące w prawo, a prawa krawędź współrzędne y rosnące w dół.} \\
\texttt{* Następujące obiekty są losowo umieszczone na siatce: 'P', 'O', 'T', 'S', 'D', 'A', 'M'.} \\
\\ 
\\ 
\texttt{Drugi gracz widzi wariant planszy, na której obiekty są umieszczone w różnych losowych miejscach. Nie można zobaczyć planszy drugiego gracza, a on nie może zobaczyć Państwa planszy.} \\
\\ 
\\ 
\texttt{**Cel:**} \\
\\ 
\texttt{Obaj gracze muszą przesuwać obiekty na swoim tle tak, aby identyczne obiekty znalazły się na tych samych współrzędnych. Muszą Państwo komunikować się z drugim graczem, aby uzgodnić wspólną konfigurację docelową.} \\
\\ 
\\ 
\texttt{**Zasady:**} \\
\\ 
\texttt{* W każdej turze można wysłać dokładnie jedno z dwóch poniższych poleceń:} \\
\texttt{ 1. `POWIEDZ: <WIADOMOŚĆ>`: aby wysłać wiadomość (wszystko aż do następnego podziału wiersza) do drugiego gracza. Przekażę ją Państwa partnerowi.} \\
\texttt{2. `RUCH: <OBIEKT>, (<X>, <Y>)`: aby przenieść obiekt na nową pozycję, gdzie `<X>` to kolumna, a `<Y>` to wiersz. Poinformuję Państwa, czy przeniesienie się powiodło, czy nie.} \\
\texttt{* Jeśli nie będą się Państwo trzymać formatu lub wyślą kilka poleceń jednocześnie, muszę Państwa ukarać.} \\
\texttt{* Jeśli obaj gracze zgromadzą więcej niż 16 kar, obaj przegrają grę.} \\
\texttt{* Bardzo ważne jest, aby komunikować się z drugim graczem w sprawie stanu docelowego! *Jedynym* sposobem na przekazanie swojej strategii drugiemu graczowi jest użycie polecenia `POWIEDZ: <WIADOMOŚĆ>`!} \\
\\ 
\\ 
\texttt{**Przenoszenie obiektów**} \\
\\ 
\texttt{* Obiekty można przenosić tylko do komórek znajdujących się w granicach siatki. Komórka docelowa musi być pusta, tzn. musi zawierać tylko symbol '◌'.} \\
\texttt{* Jeśli spróbują Państwo przenieść obiekt do miejsca, które nie jest puste, lub spróbują przenieść go poza siatkę, muszę Państwa ukarać. Otrzymają Państwo kolejną próbę.} \\
\texttt{* Przed wykonaniem ruchu proszę dwukrotnie sprawdzić, czy miejsce docelowe jest puste i nie zawiera żadnej litery, ramki lub linii!} \\
\\ 
\\ 
\texttt{**Koniec gry**} \\
\\ 
\texttt{Jeśli uważają Państwo, że osiągnęli cel wyrównania wszystkich obiektów, mogą Państwo poprosić drugiego gracza o zakończenie gry, wysyłając `POWIEDZ: zakończono?`. Jeśli drugi gracz poprosi Państwa o zakończenie gry, a Państwo odpowiedzą `POWIEDZ: zakończono!`, gra się zakończy.} \\
\\ 
\texttt{Obaj gracze wygrywają, jeśli gra zakończy się w ciągu 28 rund, gdzie jedna runda jest definiowana jako wysłanie ważnego polecenia przez dwóch graczy.} \\
\\ 
\texttt{**Punktacja:**} \\
\\ 
\texttt{Im bliżej siebie znajdują się identyczne obiekty na obu planszach, tym więcej punktów otrzymują Państwo. Kary zmniejszają liczbę punktów. Czy uda się Państwu pobić rekord?} \\
\\ 
\texttt{Drugi gracz rozpoczął grę wysyłając następującą wiadomość:} \\
\\ 
\texttt{"Gotowy do rozpoczęcia gry!"} \\
\\ 
\texttt{Jakie jest twoje pierwsze polecenie?} \\
            }
        }
    }
    & \\ \\

    \theutterance \stepcounter{utterance}  
    & & & \multicolumn{4}{p{0.6\linewidth}}{
        \cellcolor[rgb]{0.9,0.9,0.9}{
            \makecell[{{p{\linewidth}}}]{
                \texttt{\tiny{[GM$\langle$P2]}}
                \texttt{RUCH: P, (2, 1)} \\
            }
        }
    }
     \\ \\

    \theutterance \stepcounter{utterance}  
    & & & \multicolumn{2}{p{0.3\linewidth}}{
        \cellcolor[rgb]{0.9,0.9,0.9}{
            \makecell[{{p{\linewidth}}}]{
                \texttt{\tiny{[GM$|$GM]}}
                \texttt{RUCH: P, (2, 1)} \\
            }
        }
    }
    & & \\ \\

    \theutterance \stepcounter{utterance}  
    & & & \multicolumn{2}{p{0.3\linewidth}}{
        \cellcolor[rgb]{0.9,0.9,0.9}{
            \makecell[{{p{\linewidth}}}]{
                \texttt{\tiny{[GM$|$GM]}}
                \texttt{P moved to (2, 1)} \\
            }
        }
    }
    & & \\ \\

    \theutterance \stepcounter{utterance}  
    & & \multicolumn{4}{p{0.6\linewidth}}{
        \cellcolor[rgb]{0.9,0.9,0.9}{
            \makecell[{{p{\linewidth}}}]{
                \texttt{\tiny{[P1$\langle$GM]}}
                \texttt{Twoja wiadomość została przekazana drugiemu graczowi.} \\
\\ 
\\ 
\texttt{Obecnie grasz rundę 1 z maksymalnie 28.} \\
\texttt{Łącznie zgromadziliście 0 z 16 kar.} \\
\texttt{Drugi gracz przesunął obiekt na swojej siatce.} \\
\\ 
\\ 
\texttt{Jakie jest twoje następne polecenie?} \\
            }
        }
    }
    & \\ \\

    \theutterance \stepcounter{utterance}  
    & \multicolumn{4}{p{0.6\linewidth}}{
        \cellcolor[rgb]{0.9,0.9,0.9}{
            \makecell[{{p{\linewidth}}}]{
                \texttt{\tiny{[P1$\rangle$GM]}}
                \texttt{RUCH: P, (1, 2)} \\
            }
        }
    }
    & & \\ \\

    \theutterance \stepcounter{utterance}  
    & & & \multicolumn{2}{p{0.3\linewidth}}{
        \cellcolor[rgb]{0.9,0.9,0.9}{
            \makecell[{{p{\linewidth}}}]{
                \texttt{\tiny{[GM$|$GM]}}
                \texttt{RUCH: P, (1, 2)} \\
            }
        }
    }
    & & \\ \\

    \theutterance \stepcounter{utterance}  
    & & & \multicolumn{2}{p{0.3\linewidth}}{
        \cellcolor[rgb]{0.9,0.9,0.9}{
            \makecell[{{p{\linewidth}}}]{
                \texttt{\tiny{[GM$|$GM]}}
                \texttt{P moved to (1, 2)} \\
            }
        }
    }
    & & \\ \\

    \theutterance \stepcounter{utterance}  
    & & \multicolumn{4}{p{0.6\linewidth}}{
        \cellcolor[rgb]{0.9,0.9,0.9}{
            \makecell[{{p{\linewidth}}}]{
                \texttt{\tiny{[GM$\rangle$P2]}}
                \texttt{Pomyślnie przeniesiono 'P' do (2,1). Twoja zaktualizowana siatka wygląda następująco:} \\
\texttt{```} \\
\texttt{ 1234567} \\
\texttt{╔════╤══╗} \\
\texttt{║◌P◌◌│O◌║ 1} \\
\texttt{║◌◌M◌├──╢ 2} \\
\texttt{║◌S◌◌│T◌║ 3} \\
\texttt{║◌┌──┼──╢ 4} \\
\texttt{║A│◌◌│◌◌║ 5} \\
\texttt{╟─┴┬─┤◌D║ 6} \\
\texttt{║◌◌│◌│◌◌║ 7} \\
\texttt{╚══╧═╧══╝} \\
\texttt{```} \\
\texttt{Obecnie grasz rundę 2 z maksymalnie 28.} \\
\texttt{Łącznie zgromadziliście 0 z 16 kar.} \\
\texttt{Drugi gracz przesunął obiekt na swojej siatce.} \\
\\ 
\\ 
\texttt{Jakie jest twoje następne polecenie?} \\
            }
        }
    }
    & \\ \\

    \theutterance \stepcounter{utterance}  
    & & & \multicolumn{4}{p{0.6\linewidth}}{
        \cellcolor[rgb]{0.9,0.9,0.9}{
            \makecell[{{p{\linewidth}}}]{
                \texttt{\tiny{[GM$\langle$P2]}}
                \texttt{POWIEDZ: Proszę o aktualną sytuację docelową dla obiektów 'O', 'T', 'S', 'D', 'A', 'M'.} \\
            }
        }
    }
     \\ \\

    \theutterance \stepcounter{utterance}  
    & & & \multicolumn{2}{p{0.3\linewidth}}{
        \cellcolor[rgb]{0.9,0.9,0.9}{
            \makecell[{{p{\linewidth}}}]{
                \texttt{\tiny{[GM$|$GM]}}
                \texttt{POWIEDZ: Proszę o aktualną sytuację docelową dla obiektów 'O', 'T', 'S', 'D', 'A', 'M'.} \\
            }
        }
    }
    & & \\ \\

    \theutterance \stepcounter{utterance}  
    & & \multicolumn{4}{p{0.6\linewidth}}{
        \cellcolor[rgb]{0.9,0.9,0.9}{
            \makecell[{{p{\linewidth}}}]{
                \texttt{\tiny{[P1$\langle$GM]}}
                \texttt{Pomyślnie przeniesiono 'P' do (1,2). Twoja zaktualizowana siatka wygląda następująco:} \\
\texttt{```} \\
\texttt{ 1234567} \\
\texttt{╔════╤══╗} \\
\texttt{║◌◌◌O│◌M║ 1} \\
\texttt{║P◌T◌├──╢ 2} \\
\texttt{║◌◌◌◌│◌◌║ 3} \\
\texttt{║S┌──┼──╢ 4} \\
\texttt{║◌│◌A│◌◌║ 5} \\
\texttt{╟─┴┬─┤◌◌║ 6} \\
\texttt{║◌◌│◌│D◌║ 7} \\
\texttt{╚══╧═╧══╝} \\
\texttt{```} \\
\texttt{Obecnie grasz rundę 2 z maksymalnie 28.} \\
\texttt{Łącznie zgromadziliście 0 z 16 kar.} \\
\texttt{Drugi gracz wysłał tę wiadomość:} \\
\\ 
\\ 
\texttt{"Proszę o aktualną sytuację docelową dla obiektów 'O', 'T', 'S', 'D', 'A', 'M'."} \\
\\ 
\\ 
\texttt{Jakie jest twoje następne polecenie?} \\
            }
        }
    }
    & \\ \\

    \theutterance \stepcounter{utterance}  
    & \multicolumn{4}{p{0.6\linewidth}}{
        \cellcolor[rgb]{0.9,0.9,0.9}{
            \makecell[{{p{\linewidth}}}]{
                \texttt{\tiny{[P1$\rangle$GM]}}
                \texttt{POWIEDZ: Obecnie nie mam ustalonej konfiguracji docelowej. Proszę pomóc ustalić, gdzie powinny znaleźć się obiekty 'O', 'T', 'S', 'D', 'A', 'M'. Proponuję następujące lokalizacje docelowe:} \\
\texttt{{-} 'O' na (2,1)} \\
\texttt{{-} 'T' na (2,3)} \\
\texttt{{-} 'S' na (4,1)} \\
\texttt{{-} 'D' na (7,5)} \\
\texttt{{-} 'A' na (5,3)} \\
\texttt{{-} 'M' na (1,6)} \\
            }
        }
    }
    & & \\ \\

    \theutterance \stepcounter{utterance}  
    & & & \multicolumn{2}{p{0.3\linewidth}}{
        \cellcolor[rgb]{0.9,0.9,0.9}{
            \makecell[{{p{\linewidth}}}]{
                \texttt{\tiny{[GM$|$GM]}}
                \texttt{POWIEDZ: Obecnie nie mam ustalonej konfiguracji docelowej. Proszę pomóc ustalić, gdzie powinny znaleźć się obiekty 'O', 'T', 'S', 'D', 'A', 'M'. Proponuję następujące lokalizacje docelowe:} \\
\texttt{{-} 'O' na (2,1)} \\
\texttt{{-} 'T' na (2,3)} \\
\texttt{{-} 'S' na (4,1)} \\
\texttt{{-} 'D' na (7,5)} \\
\texttt{{-} 'A' na (5,3)} \\
\texttt{{-} 'M' na (1,6)} \\
            }
        }
    }
    & & \\ \\

    \theutterance \stepcounter{utterance}  
    & & & \multicolumn{2}{p{0.3\linewidth}}{
        \cellcolor[rgb]{0.9,0.9,0.9}{
            \makecell[{{p{\linewidth}}}]{
                \texttt{\tiny{[GM$|$GM]}}
                \texttt{Invalid format: tail is not empty: '} \\
\texttt{{-} 'O' na (2,1)} \\
\texttt{{-} 'T' na (2,3)} \\
\texttt{{-} 'S' na (4,1)} \\
\texttt{{-} 'D' na (7,5)} \\
\texttt{{-} 'A' na (5,3)} \\
\texttt{{-} 'M' na (1,6)} \\
            }
        }
    }
    & & \\ \\

    \theutterance \stepcounter{utterance}  
    & & \multicolumn{4}{p{0.6\linewidth}}{
        \cellcolor[rgb]{0.9,0.9,0.9}{
            \makecell[{{p{\linewidth}}}]{
                \texttt{\tiny{[P1$\langle$GM]}}
                \texttt{Kara: Twoja wiadomość nie może zawierać niczego po poleceniu!} \\
\\ 
\texttt{Upewnij się, że twoja odpowiedź zawiera tylko `POWIEDZ: <WIADOMOŚĆ>` lub `RUCH: <OBIEKT>, (<X>, <Y>)` i nic więcej!} \\
\\ 
\texttt{Łącznie zgromadziliście 1 z 16 kar.} \\
\texttt{Spróbuj ponownie!} \\
            }
        }
    }
    & \\ \\

    \theutterance \stepcounter{utterance}  
    & \multicolumn{4}{p{0.6\linewidth}}{
        \cellcolor[rgb]{0.9,0.9,0.9}{
            \makecell[{{p{\linewidth}}}]{
                \texttt{\tiny{[P1$\rangle$GM]}}
                \texttt{POWIEDZ: Gotowy do kontynuowania.} \\
            }
        }
    }
    & & \\ \\

    \theutterance \stepcounter{utterance}  
    & & & \multicolumn{2}{p{0.3\linewidth}}{
        \cellcolor[rgb]{0.9,0.9,0.9}{
            \makecell[{{p{\linewidth}}}]{
                \texttt{\tiny{[GM$|$GM]}}
                \texttt{POWIEDZ: Gotowy do kontynuowania.} \\
            }
        }
    }
    & & \\ \\

    \theutterance \stepcounter{utterance}  
    & & \multicolumn{4}{p{0.6\linewidth}}{
        \cellcolor[rgb]{0.9,0.9,0.9}{
            \makecell[{{p{\linewidth}}}]{
                \texttt{\tiny{[GM$\rangle$P2]}}
                \texttt{Twoja wiadomość została przekazana drugiemu graczowi.} \\
\\ 
\\ 
\texttt{Obecnie grasz rundę 3 z maksymalnie 28.} \\
\texttt{Łącznie zgromadziliście 1 z 16 kar.} \\
\texttt{Drugi gracz wysłał tę wiadomość:} \\
\\ 
\\ 
\texttt{"Gotowy do kontynuowania."} \\
\\ 
\\ 
\texttt{Jakie jest twoje następne polecenie?} \\
            }
        }
    }
    & \\ \\

    \theutterance \stepcounter{utterance}  
    & & & \multicolumn{4}{p{0.6\linewidth}}{
        \cellcolor[rgb]{0.9,0.9,0.9}{
            \makecell[{{p{\linewidth}}}]{
                \texttt{\tiny{[GM$\langle$P2]}}
                \texttt{RUCH: O, (1, 4)} \\
            }
        }
    }
     \\ \\

    \theutterance \stepcounter{utterance}  
    & & & \multicolumn{2}{p{0.3\linewidth}}{
        \cellcolor[rgb]{0.9,0.9,0.9}{
            \makecell[{{p{\linewidth}}}]{
                \texttt{\tiny{[GM$|$GM]}}
                \texttt{RUCH: O, (1, 4)} \\
            }
        }
    }
    & & \\ \\

    \theutterance \stepcounter{utterance}  
    & & & \multicolumn{2}{p{0.3\linewidth}}{
        \cellcolor[rgb]{0.9,0.9,0.9}{
            \makecell[{{p{\linewidth}}}]{
                \texttt{\tiny{[GM$|$GM]}}
                \texttt{O moved to (1, 4)} \\
            }
        }
    }
    & & \\ \\

    \theutterance \stepcounter{utterance}  
    & & \multicolumn{4}{p{0.6\linewidth}}{
        \cellcolor[rgb]{0.9,0.9,0.9}{
            \makecell[{{p{\linewidth}}}]{
                \texttt{\tiny{[P1$\langle$GM]}}
                \texttt{Twoja wiadomość została przekazana drugiemu graczowi.} \\
\\ 
\\ 
\texttt{Obecnie grasz rundę 3 z maksymalnie 28.} \\
\texttt{Łącznie zgromadziliście 1 z 16 kar.} \\
\texttt{Drugi gracz przesunął obiekt na swojej siatce.} \\
\\ 
\\ 
\texttt{Jakie jest twoje następne polecenie?} \\
            }
        }
    }
    & \\ \\

    \theutterance \stepcounter{utterance}  
    & \multicolumn{4}{p{0.6\linewidth}}{
        \cellcolor[rgb]{0.9,0.9,0.9}{
            \makecell[{{p{\linewidth}}}]{
                \texttt{\tiny{[P1$\rangle$GM]}}
                \texttt{RUCH: T, (2, 3)} \\
            }
        }
    }
    & & \\ \\

    \theutterance \stepcounter{utterance}  
    & & & \multicolumn{2}{p{0.3\linewidth}}{
        \cellcolor[rgb]{0.9,0.9,0.9}{
            \makecell[{{p{\linewidth}}}]{
                \texttt{\tiny{[GM$|$GM]}}
                \texttt{RUCH: T, (2, 3)} \\
            }
        }
    }
    & & \\ \\

    \theutterance \stepcounter{utterance}  
    & & & \multicolumn{2}{p{0.3\linewidth}}{
        \cellcolor[rgb]{0.9,0.9,0.9}{
            \makecell[{{p{\linewidth}}}]{
                \texttt{\tiny{[GM$|$GM]}}
                \texttt{T moved to (2, 3)} \\
            }
        }
    }
    & & \\ \\

    \theutterance \stepcounter{utterance}  
    & & \multicolumn{4}{p{0.6\linewidth}}{
        \cellcolor[rgb]{0.9,0.9,0.9}{
            \makecell[{{p{\linewidth}}}]{
                \texttt{\tiny{[GM$\rangle$P2]}}
                \texttt{Pomyślnie przeniesiono 'O' do (1,4). Twoja zaktualizowana siatka wygląda następująco:} \\
\texttt{```} \\
\texttt{ 1234567} \\
\texttt{╔════╤══╗} \\
\texttt{║◌P◌◌│◌◌║ 1} \\
\texttt{║◌◌M◌├──╢ 2} \\
\texttt{║◌S◌◌│T◌║ 3} \\
\texttt{║O┌──┼──╢ 4} \\
\texttt{║A│◌◌│◌◌║ 5} \\
\texttt{╟─┴┬─┤◌D║ 6} \\
\texttt{║◌◌│◌│◌◌║ 7} \\
\texttt{╚══╧═╧══╝} \\
\texttt{```} \\
\texttt{Obecnie grasz rundę 4 z maksymalnie 28.} \\
\texttt{Łącznie zgromadziliście 1 z 16 kar.} \\
\texttt{Drugi gracz przesunął obiekt na swojej siatce.} \\
\\ 
\\ 
\texttt{Jakie jest twoje następne polecenie?} \\
            }
        }
    }
    & \\ \\

    \theutterance \stepcounter{utterance}  
    & & & \multicolumn{4}{p{0.6\linewidth}}{
        \cellcolor[rgb]{0.9,0.9,0.9}{
            \makecell[{{p{\linewidth}}}]{
                \texttt{\tiny{[GM$\langle$P2]}}
                \texttt{RUCH: T, (3, 5)} \\
            }
        }
    }
     \\ \\

    \theutterance \stepcounter{utterance}  
    & & & \multicolumn{2}{p{0.3\linewidth}}{
        \cellcolor[rgb]{0.9,0.9,0.9}{
            \makecell[{{p{\linewidth}}}]{
                \texttt{\tiny{[GM$|$GM]}}
                \texttt{RUCH: T, (3, 5)} \\
            }
        }
    }
    & & \\ \\

    \theutterance \stepcounter{utterance}  
    & & & \multicolumn{2}{p{0.3\linewidth}}{
        \cellcolor[rgb]{0.9,0.9,0.9}{
            \makecell[{{p{\linewidth}}}]{
                \texttt{\tiny{[GM$|$GM]}}
                \texttt{T moved to (3, 5)} \\
            }
        }
    }
    & & \\ \\

    \theutterance \stepcounter{utterance}  
    & & \multicolumn{4}{p{0.6\linewidth}}{
        \cellcolor[rgb]{0.9,0.9,0.9}{
            \makecell[{{p{\linewidth}}}]{
                \texttt{\tiny{[P1$\langle$GM]}}
                \texttt{Pomyślnie przeniesiono 'T' do (2,3). Twoja zaktualizowana siatka wygląda następująco:} \\
\texttt{```} \\
\texttt{ 1234567} \\
\texttt{╔════╤══╗} \\
\texttt{║◌◌◌O│◌M║ 1} \\
\texttt{║P◌◌◌├──╢ 2} \\
\texttt{║◌T◌◌│◌◌║ 3} \\
\texttt{║S┌──┼──╢ 4} \\
\texttt{║◌│◌A│◌◌║ 5} \\
\texttt{╟─┴┬─┤◌◌║ 6} \\
\texttt{║◌◌│◌│D◌║ 7} \\
\texttt{╚══╧═╧══╝} \\
\texttt{```} \\
\texttt{Obecnie grasz rundę 4 z maksymalnie 28.} \\
\texttt{Łącznie zgromadziliście 1 z 16 kar.} \\
\texttt{Drugi gracz przesunął obiekt na swojej siatce.} \\
\\ 
\\ 
\texttt{Jakie jest twoje następne polecenie?} \\
            }
        }
    }
    & \\ \\

    \theutterance \stepcounter{utterance}  
    & \multicolumn{4}{p{0.6\linewidth}}{
        \cellcolor[rgb]{0.9,0.9,0.9}{
            \makecell[{{p{\linewidth}}}]{
                \texttt{\tiny{[P1$\rangle$GM]}}
                \texttt{RUCH: S, (4, 1)} \\
            }
        }
    }
    & & \\ \\

    \theutterance \stepcounter{utterance}  
    & & & \multicolumn{2}{p{0.3\linewidth}}{
        \cellcolor[rgb]{0.9,0.9,0.9}{
            \makecell[{{p{\linewidth}}}]{
                \texttt{\tiny{[GM$|$GM]}}
                \texttt{RUCH: S, (4, 1)} \\
            }
        }
    }
    & & \\ \\

    \theutterance \stepcounter{utterance}  
    & & & \multicolumn{2}{p{0.3\linewidth}}{
        \cellcolor[rgb]{0.9,0.9,0.9}{
            \makecell[{{p{\linewidth}}}]{
                \texttt{\tiny{[GM$|$GM]}}
                \texttt{Kara: (4,1) nie jest puste, ale zawiera 'O'.} \\
\\ 
\texttt{Łącznie zgromadziliście 2 z 16 kar.} \\
\texttt{Spróbuj ponownie!} \\
            }
        }
    }
    & & \\ \\

    \theutterance \stepcounter{utterance}  
    & & \multicolumn{4}{p{0.6\linewidth}}{
        \cellcolor[rgb]{0.9,0.9,0.9}{
            \makecell[{{p{\linewidth}}}]{
                \texttt{\tiny{[P1$\langle$GM]}}
                \texttt{Kara: (4,1) nie jest puste, ale zawiera 'O'.} \\
\\ 
\texttt{Łącznie zgromadziliście 2 z 16 kar.} \\
\texttt{Spróbuj ponownie!} \\
            }
        }
    }
    & \\ \\

    \theutterance \stepcounter{utterance}  
    & \multicolumn{4}{p{0.6\linewidth}}{
        \cellcolor[rgb]{0.9,0.9,0.9}{
            \makecell[{{p{\linewidth}}}]{
                \texttt{\tiny{[P1$\rangle$GM]}}
                \texttt{RUCH: S, (5, 1)} \\
            }
        }
    }
    & & \\ \\

    \theutterance \stepcounter{utterance}  
    & & & \multicolumn{2}{p{0.3\linewidth}}{
        \cellcolor[rgb]{0.9,0.9,0.9}{
            \makecell[{{p{\linewidth}}}]{
                \texttt{\tiny{[GM$|$GM]}}
                \texttt{RUCH: S, (5, 1)} \\
            }
        }
    }
    & & \\ \\

    \theutterance \stepcounter{utterance}  
    & & & \multicolumn{2}{p{0.3\linewidth}}{
        \cellcolor[rgb]{0.9,0.9,0.9}{
            \makecell[{{p{\linewidth}}}]{
                \texttt{\tiny{[GM$|$GM]}}
                \texttt{Kara: (5,1) nie jest puste, ale zawiera '│'.} \\
\\ 
\texttt{Łącznie zgromadziliście 3 z 16 kar.} \\
\texttt{Spróbuj ponownie!} \\
            }
        }
    }
    & & \\ \\

    \theutterance \stepcounter{utterance}  
    & & \multicolumn{4}{p{0.6\linewidth}}{
        \cellcolor[rgb]{0.9,0.9,0.9}{
            \makecell[{{p{\linewidth}}}]{
                \texttt{\tiny{[P1$\langle$GM]}}
                \texttt{Kara: (5,1) nie jest puste, ale zawiera '│'.} \\
\\ 
\texttt{Łącznie zgromadziliście 3 z 16 kar.} \\
\texttt{Spróbuj ponownie!} \\
            }
        }
    }
    & \\ \\

    \theutterance \stepcounter{utterance}  
    & \multicolumn{4}{p{0.6\linewidth}}{
        \cellcolor[rgb]{0.9,0.9,0.9}{
            \makecell[{{p{\linewidth}}}]{
                \texttt{\tiny{[P1$\rangle$GM]}}
                \texttt{RUCH: S, (6, 1)} \\
            }
        }
    }
    & & \\ \\

    \theutterance \stepcounter{utterance}  
    & & & \multicolumn{2}{p{0.3\linewidth}}{
        \cellcolor[rgb]{0.9,0.9,0.9}{
            \makecell[{{p{\linewidth}}}]{
                \texttt{\tiny{[GM$|$GM]}}
                \texttt{RUCH: S, (6, 1)} \\
            }
        }
    }
    & & \\ \\

    \theutterance \stepcounter{utterance}  
    & & & \multicolumn{2}{p{0.3\linewidth}}{
        \cellcolor[rgb]{0.9,0.9,0.9}{
            \makecell[{{p{\linewidth}}}]{
                \texttt{\tiny{[GM$|$GM]}}
                \texttt{S moved to (6, 1)} \\
            }
        }
    }
    & & \\ \\

    \theutterance \stepcounter{utterance}  
    & & \multicolumn{4}{p{0.6\linewidth}}{
        \cellcolor[rgb]{0.9,0.9,0.9}{
            \makecell[{{p{\linewidth}}}]{
                \texttt{\tiny{[GM$\rangle$P2]}}
                \texttt{Pomyślnie przeniesiono 'T' do (3,5). Twoja zaktualizowana siatka wygląda następująco:} \\
\texttt{```} \\
\texttt{ 1234567} \\
\texttt{╔════╤══╗} \\
\texttt{║◌P◌◌│◌◌║ 1} \\
\texttt{║◌◌M◌├──╢ 2} \\
\texttt{║◌S◌◌│◌◌║ 3} \\
\texttt{║O┌──┼──╢ 4} \\
\texttt{║A│T◌│◌◌║ 5} \\
\texttt{╟─┴┬─┤◌D║ 6} \\
\texttt{║◌◌│◌│◌◌║ 7} \\
\texttt{╚══╧═╧══╝} \\
\texttt{```} \\
\texttt{Obecnie grasz rundę 5 z maksymalnie 28.} \\
\texttt{Łącznie zgromadziliście 3 z 16 kar.} \\
\texttt{Drugi gracz przesunął obiekt na swojej siatce.} \\
\\ 
\\ 
\texttt{Jakie jest twoje następne polecenie?} \\
            }
        }
    }
    & \\ \\

    \theutterance \stepcounter{utterance}  
    & & & \multicolumn{4}{p{0.6\linewidth}}{
        \cellcolor[rgb]{0.9,0.9,0.9}{
            \makecell[{{p{\linewidth}}}]{
                \texttt{\tiny{[GM$\langle$P2]}}
                \texttt{RUCH: S, (2, 3)} \\
            }
        }
    }
     \\ \\

    \theutterance \stepcounter{utterance}  
    & & & \multicolumn{2}{p{0.3\linewidth}}{
        \cellcolor[rgb]{0.9,0.9,0.9}{
            \makecell[{{p{\linewidth}}}]{
                \texttt{\tiny{[GM$|$GM]}}
                \texttt{RUCH: S, (2, 3)} \\
            }
        }
    }
    & & \\ \\

    \theutterance \stepcounter{utterance}  
    & & & \multicolumn{2}{p{0.3\linewidth}}{
        \cellcolor[rgb]{0.9,0.9,0.9}{
            \makecell[{{p{\linewidth}}}]{
                \texttt{\tiny{[GM$|$GM]}}
                \texttt{Kara: (2,3) nie jest puste, ale zawiera 'S'.} \\
\\ 
\texttt{Łącznie zgromadziliście 4 z 16 kar.} \\
\texttt{Spróbuj ponownie!} \\
            }
        }
    }
    & & \\ \\

    \theutterance \stepcounter{utterance}  
    & & \multicolumn{4}{p{0.6\linewidth}}{
        \cellcolor[rgb]{0.9,0.9,0.9}{
            \makecell[{{p{\linewidth}}}]{
                \texttt{\tiny{[GM$\rangle$P2]}}
                \texttt{Kara: (2,3) nie jest puste, ale zawiera 'S'.} \\
\\ 
\texttt{Łącznie zgromadziliście 4 z 16 kar.} \\
\texttt{Spróbuj ponownie!} \\
            }
        }
    }
    & \\ \\

    \theutterance \stepcounter{utterance}  
    & & & \multicolumn{4}{p{0.6\linewidth}}{
        \cellcolor[rgb]{0.9,0.9,0.9}{
            \makecell[{{p{\linewidth}}}]{
                \texttt{\tiny{[GM$\langle$P2]}}
                \texttt{RUCH: S, (2, 4)} \\
            }
        }
    }
     \\ \\

    \theutterance \stepcounter{utterance}  
    & & & \multicolumn{2}{p{0.3\linewidth}}{
        \cellcolor[rgb]{0.9,0.9,0.9}{
            \makecell[{{p{\linewidth}}}]{
                \texttt{\tiny{[GM$|$GM]}}
                \texttt{RUCH: S, (2, 4)} \\
            }
        }
    }
    & & \\ \\

    \theutterance \stepcounter{utterance}  
    & & & \multicolumn{2}{p{0.3\linewidth}}{
        \cellcolor[rgb]{0.9,0.9,0.9}{
            \makecell[{{p{\linewidth}}}]{
                \texttt{\tiny{[GM$|$GM]}}
                \texttt{Kara: (2,4) nie jest puste, ale zawiera '┌'.} \\
\\ 
\texttt{Łącznie zgromadziliście 5 z 16 kar.} \\
\texttt{Spróbuj ponownie!} \\
            }
        }
    }
    & & \\ \\

    \theutterance \stepcounter{utterance}  
    & & \multicolumn{4}{p{0.6\linewidth}}{
        \cellcolor[rgb]{0.9,0.9,0.9}{
            \makecell[{{p{\linewidth}}}]{
                \texttt{\tiny{[GM$\rangle$P2]}}
                \texttt{Kara: (2,4) nie jest puste, ale zawiera '┌'.} \\
\\ 
\texttt{Łącznie zgromadziliście 5 z 16 kar.} \\
\texttt{Spróbuj ponownie!} \\
            }
        }
    }
    & \\ \\

    \theutterance \stepcounter{utterance}  
    & & & \multicolumn{4}{p{0.6\linewidth}}{
        \cellcolor[rgb]{0.9,0.9,0.9}{
            \makecell[{{p{\linewidth}}}]{
                \texttt{\tiny{[GM$\langle$P2]}}
                \texttt{RUCH: S, (3, 3)} \\
            }
        }
    }
     \\ \\

    \theutterance \stepcounter{utterance}  
    & & & \multicolumn{2}{p{0.3\linewidth}}{
        \cellcolor[rgb]{0.9,0.9,0.9}{
            \makecell[{{p{\linewidth}}}]{
                \texttt{\tiny{[GM$|$GM]}}
                \texttt{RUCH: S, (3, 3)} \\
            }
        }
    }
    & & \\ \\

    \theutterance \stepcounter{utterance}  
    & & & \multicolumn{2}{p{0.3\linewidth}}{
        \cellcolor[rgb]{0.9,0.9,0.9}{
            \makecell[{{p{\linewidth}}}]{
                \texttt{\tiny{[GM$|$GM]}}
                \texttt{S moved to (3, 3)} \\
            }
        }
    }
    & & \\ \\

    \theutterance \stepcounter{utterance}  
    & & \multicolumn{4}{p{0.6\linewidth}}{
        \cellcolor[rgb]{0.9,0.9,0.9}{
            \makecell[{{p{\linewidth}}}]{
                \texttt{\tiny{[P1$\langle$GM]}}
                \texttt{Pomyślnie przeniesiono 'S' do (6,1). Twoja zaktualizowana siatka wygląda następująco:} \\
\texttt{```} \\
\texttt{ 1234567} \\
\texttt{╔════╤══╗} \\
\texttt{║◌◌◌O│SM║ 1} \\
\texttt{║P◌◌◌├──╢ 2} \\
\texttt{║◌T◌◌│◌◌║ 3} \\
\texttt{║◌┌──┼──╢ 4} \\
\texttt{║◌│◌A│◌◌║ 5} \\
\texttt{╟─┴┬─┤◌◌║ 6} \\
\texttt{║◌◌│◌│D◌║ 7} \\
\texttt{╚══╧═╧══╝} \\
\texttt{```} \\
\texttt{Obecnie grasz rundę 5 z maksymalnie 28.} \\
\texttt{Łącznie zgromadziliście 5 z 16 kar.} \\
\texttt{Drugi gracz przesunął obiekt na swojej siatce.} \\
\\ 
\\ 
\texttt{Jakie jest twoje następne polecenie?} \\
            }
        }
    }
    & \\ \\

    \theutterance \stepcounter{utterance}  
    & \multicolumn{4}{p{0.6\linewidth}}{
        \cellcolor[rgb]{0.9,0.9,0.9}{
            \makecell[{{p{\linewidth}}}]{
                \texttt{\tiny{[P1$\rangle$GM]}}
                \texttt{RUCH: A, (5, 3)} \\
            }
        }
    }
    & & \\ \\

    \theutterance \stepcounter{utterance}  
    & & & \multicolumn{2}{p{0.3\linewidth}}{
        \cellcolor[rgb]{0.9,0.9,0.9}{
            \makecell[{{p{\linewidth}}}]{
                \texttt{\tiny{[GM$|$GM]}}
                \texttt{RUCH: A, (5, 3)} \\
            }
        }
    }
    & & \\ \\

    \theutterance \stepcounter{utterance}  
    & & & \multicolumn{2}{p{0.3\linewidth}}{
        \cellcolor[rgb]{0.9,0.9,0.9}{
            \makecell[{{p{\linewidth}}}]{
                \texttt{\tiny{[GM$|$GM]}}
                \texttt{Kara: (5,3) nie jest puste, ale zawiera '│'.} \\
\\ 
\texttt{Łącznie zgromadziliście 6 z 16 kar.} \\
\texttt{Spróbuj ponownie!} \\
            }
        }
    }
    & & \\ \\

    \theutterance \stepcounter{utterance}  
    & & \multicolumn{4}{p{0.6\linewidth}}{
        \cellcolor[rgb]{0.9,0.9,0.9}{
            \makecell[{{p{\linewidth}}}]{
                \texttt{\tiny{[P1$\langle$GM]}}
                \texttt{Kara: (5,3) nie jest puste, ale zawiera '│'.} \\
\\ 
\texttt{Łącznie zgromadziliście 6 z 16 kar.} \\
\texttt{Spróbuj ponownie!} \\
            }
        }
    }
    & \\ \\

    \theutterance \stepcounter{utterance}  
    & \multicolumn{4}{p{0.6\linewidth}}{
        \cellcolor[rgb]{0.9,0.9,0.9}{
            \makecell[{{p{\linewidth}}}]{
                \texttt{\tiny{[P1$\rangle$GM]}}
                \texttt{RUCH: A, (4, 3)} \\
            }
        }
    }
    & & \\ \\

    \theutterance \stepcounter{utterance}  
    & & & \multicolumn{2}{p{0.3\linewidth}}{
        \cellcolor[rgb]{0.9,0.9,0.9}{
            \makecell[{{p{\linewidth}}}]{
                \texttt{\tiny{[GM$|$GM]}}
                \texttt{RUCH: A, (4, 3)} \\
            }
        }
    }
    & & \\ \\

    \theutterance \stepcounter{utterance}  
    & & & \multicolumn{2}{p{0.3\linewidth}}{
        \cellcolor[rgb]{0.9,0.9,0.9}{
            \makecell[{{p{\linewidth}}}]{
                \texttt{\tiny{[GM$|$GM]}}
                \texttt{A moved to (4, 3)} \\
            }
        }
    }
    & & \\ \\

    \theutterance \stepcounter{utterance}  
    & & \multicolumn{4}{p{0.6\linewidth}}{
        \cellcolor[rgb]{0.9,0.9,0.9}{
            \makecell[{{p{\linewidth}}}]{
                \texttt{\tiny{[GM$\rangle$P2]}}
                \texttt{Pomyślnie przeniesiono 'S' do (3,3). Twoja zaktualizowana siatka wygląda następująco:} \\
\texttt{```} \\
\texttt{ 1234567} \\
\texttt{╔════╤══╗} \\
\texttt{║◌P◌◌│◌◌║ 1} \\
\texttt{║◌◌M◌├──╢ 2} \\
\texttt{║◌◌S◌│◌◌║ 3} \\
\texttt{║O┌──┼──╢ 4} \\
\texttt{║A│T◌│◌◌║ 5} \\
\texttt{╟─┴┬─┤◌D║ 6} \\
\texttt{║◌◌│◌│◌◌║ 7} \\
\texttt{╚══╧═╧══╝} \\
\texttt{```} \\
\texttt{Obecnie grasz rundę 6 z maksymalnie 28.} \\
\texttt{Łącznie zgromadziliście 6 z 16 kar.} \\
\texttt{Drugi gracz przesunął obiekt na swojej siatce.} \\
\\ 
\\ 
\texttt{Jakie jest twoje następne polecenie?} \\
            }
        }
    }
    & \\ \\

    \theutterance \stepcounter{utterance}  
    & & & \multicolumn{4}{p{0.6\linewidth}}{
        \cellcolor[rgb]{0.9,0.9,0.9}{
            \makecell[{{p{\linewidth}}}]{
                \texttt{\tiny{[GM$\langle$P2]}}
                \texttt{RUCH: D, (6, 5)} \\
            }
        }
    }
     \\ \\

    \theutterance \stepcounter{utterance}  
    & & & \multicolumn{2}{p{0.3\linewidth}}{
        \cellcolor[rgb]{0.9,0.9,0.9}{
            \makecell[{{p{\linewidth}}}]{
                \texttt{\tiny{[GM$|$GM]}}
                \texttt{RUCH: D, (6, 5)} \\
            }
        }
    }
    & & \\ \\

    \theutterance \stepcounter{utterance}  
    & & & \multicolumn{2}{p{0.3\linewidth}}{
        \cellcolor[rgb]{0.9,0.9,0.9}{
            \makecell[{{p{\linewidth}}}]{
                \texttt{\tiny{[GM$|$GM]}}
                \texttt{D moved to (6, 5)} \\
            }
        }
    }
    & & \\ \\

    \theutterance \stepcounter{utterance}  
    & & \multicolumn{4}{p{0.6\linewidth}}{
        \cellcolor[rgb]{0.9,0.9,0.9}{
            \makecell[{{p{\linewidth}}}]{
                \texttt{\tiny{[P1$\langle$GM]}}
                \texttt{Pomyślnie przeniesiono 'A' do (4,3). Twoja zaktualizowana siatka wygląda następująco:} \\
\texttt{```} \\
\texttt{ 1234567} \\
\texttt{╔════╤══╗} \\
\texttt{║◌◌◌O│SM║ 1} \\
\texttt{║P◌◌◌├──╢ 2} \\
\texttt{║◌T◌A│◌◌║ 3} \\
\texttt{║◌┌──┼──╢ 4} \\
\texttt{║◌│◌◌│◌◌║ 5} \\
\texttt{╟─┴┬─┤◌◌║ 6} \\
\texttt{║◌◌│◌│D◌║ 7} \\
\texttt{╚══╧═╧══╝} \\
\texttt{```} \\
\texttt{Obecnie grasz rundę 6 z maksymalnie 28.} \\
\texttt{Łącznie zgromadziliście 6 z 16 kar.} \\
\texttt{Drugi gracz przesunął obiekt na swojej siatce.} \\
\\ 
\\ 
\texttt{Jakie jest twoje następne polecenie?} \\
            }
        }
    }
    & \\ \\

    \theutterance \stepcounter{utterance}  
    & \multicolumn{4}{p{0.6\linewidth}}{
        \cellcolor[rgb]{0.9,0.9,0.9}{
            \makecell[{{p{\linewidth}}}]{
                \texttt{\tiny{[P1$\rangle$GM]}}
                \texttt{RUCH: M, (1, 6)} \\
            }
        }
    }
    & & \\ \\

    \theutterance \stepcounter{utterance}  
    & & & \multicolumn{2}{p{0.3\linewidth}}{
        \cellcolor[rgb]{0.9,0.9,0.9}{
            \makecell[{{p{\linewidth}}}]{
                \texttt{\tiny{[GM$|$GM]}}
                \texttt{RUCH: M, (1, 6)} \\
            }
        }
    }
    & & \\ \\

    \theutterance \stepcounter{utterance}  
    & & & \multicolumn{2}{p{0.3\linewidth}}{
        \cellcolor[rgb]{0.9,0.9,0.9}{
            \makecell[{{p{\linewidth}}}]{
                \texttt{\tiny{[GM$|$GM]}}
                \texttt{Kara: (1,6) nie jest puste, ale zawiera '─'.} \\
\\ 
\texttt{Łącznie zgromadziliście 7 z 16 kar.} \\
\texttt{Spróbuj ponownie!} \\
            }
        }
    }
    & & \\ \\

    \theutterance \stepcounter{utterance}  
    & & \multicolumn{4}{p{0.6\linewidth}}{
        \cellcolor[rgb]{0.9,0.9,0.9}{
            \makecell[{{p{\linewidth}}}]{
                \texttt{\tiny{[P1$\langle$GM]}}
                \texttt{Kara: (1,6) nie jest puste, ale zawiera '─'.} \\
\\ 
\texttt{Łącznie zgromadziliście 7 z 16 kar.} \\
\texttt{Spróbuj ponownie!} \\
            }
        }
    }
    & \\ \\

    \theutterance \stepcounter{utterance}  
    & \multicolumn{4}{p{0.6\linewidth}}{
        \cellcolor[rgb]{0.9,0.9,0.9}{
            \makecell[{{p{\linewidth}}}]{
                \texttt{\tiny{[P1$\rangle$GM]}}
                \texttt{RUCH: M, (7, 6)} \\
            }
        }
    }
    & & \\ \\

    \theutterance \stepcounter{utterance}  
    & & & \multicolumn{2}{p{0.3\linewidth}}{
        \cellcolor[rgb]{0.9,0.9,0.9}{
            \makecell[{{p{\linewidth}}}]{
                \texttt{\tiny{[GM$|$GM]}}
                \texttt{RUCH: M, (7, 6)} \\
            }
        }
    }
    & & \\ \\

    \theutterance \stepcounter{utterance}  
    & & & \multicolumn{2}{p{0.3\linewidth}}{
        \cellcolor[rgb]{0.9,0.9,0.9}{
            \makecell[{{p{\linewidth}}}]{
                \texttt{\tiny{[GM$|$GM]}}
                \texttt{M moved to (7, 6)} \\
            }
        }
    }
    & & \\ \\

    \theutterance \stepcounter{utterance}  
    & & \multicolumn{4}{p{0.6\linewidth}}{
        \cellcolor[rgb]{0.9,0.9,0.9}{
            \makecell[{{p{\linewidth}}}]{
                \texttt{\tiny{[GM$\rangle$P2]}}
                \texttt{Pomyślnie przeniesiono 'D' do (6,5). Twoja zaktualizowana siatka wygląda następująco:} \\
\texttt{```} \\
\texttt{ 1234567} \\
\texttt{╔════╤══╗} \\
\texttt{║◌P◌◌│◌◌║ 1} \\
\texttt{║◌◌M◌├──╢ 2} \\
\texttt{║◌◌S◌│◌◌║ 3} \\
\texttt{║O┌──┼──╢ 4} \\
\texttt{║A│T◌│D◌║ 5} \\
\texttt{╟─┴┬─┤◌◌║ 6} \\
\texttt{║◌◌│◌│◌◌║ 7} \\
\texttt{╚══╧═╧══╝} \\
\texttt{```} \\
\texttt{Obecnie grasz rundę 7 z maksymalnie 28.} \\
\texttt{Łącznie zgromadziliście 7 z 16 kar.} \\
\texttt{Drugi gracz przesunął obiekt na swojej siatce.} \\
\\ 
\\ 
\texttt{Jakie jest twoje następne polecenie?} \\
            }
        }
    }
    & \\ \\

    \theutterance \stepcounter{utterance}  
    & & & \multicolumn{4}{p{0.6\linewidth}}{
        \cellcolor[rgb]{0.9,0.9,0.9}{
            \makecell[{{p{\linewidth}}}]{
                \texttt{\tiny{[GM$\langle$P2]}}
                \texttt{RUCH: A, (5, 2)} \\
            }
        }
    }
     \\ \\

    \theutterance \stepcounter{utterance}  
    & & & \multicolumn{2}{p{0.3\linewidth}}{
        \cellcolor[rgb]{0.9,0.9,0.9}{
            \makecell[{{p{\linewidth}}}]{
                \texttt{\tiny{[GM$|$GM]}}
                \texttt{RUCH: A, (5, 2)} \\
            }
        }
    }
    & & \\ \\

    \theutterance \stepcounter{utterance}  
    & & & \multicolumn{2}{p{0.3\linewidth}}{
        \cellcolor[rgb]{0.9,0.9,0.9}{
            \makecell[{{p{\linewidth}}}]{
                \texttt{\tiny{[GM$|$GM]}}
                \texttt{Kara: (5,2) nie jest puste, ale zawiera '├'.} \\
\\ 
\texttt{Łącznie zgromadziliście 8 z 16 kar.} \\
\texttt{Spróbuj ponownie!} \\
            }
        }
    }
    & & \\ \\

    \theutterance \stepcounter{utterance}  
    & & \multicolumn{4}{p{0.6\linewidth}}{
        \cellcolor[rgb]{0.9,0.9,0.9}{
            \makecell[{{p{\linewidth}}}]{
                \texttt{\tiny{[GM$\rangle$P2]}}
                \texttt{Kara: (5,2) nie jest puste, ale zawiera '├'.} \\
\\ 
\texttt{Łącznie zgromadziliście 8 z 16 kar.} \\
\texttt{Spróbuj ponownie!} \\
            }
        }
    }
    & \\ \\

    \theutterance \stepcounter{utterance}  
    & & & \multicolumn{4}{p{0.6\linewidth}}{
        \cellcolor[rgb]{0.9,0.9,0.9}{
            \makecell[{{p{\linewidth}}}]{
                \texttt{\tiny{[GM$\langle$P2]}}
                \texttt{RUCH: A, (4, 2)} \\
            }
        }
    }
     \\ \\

    \theutterance \stepcounter{utterance}  
    & & & \multicolumn{2}{p{0.3\linewidth}}{
        \cellcolor[rgb]{0.9,0.9,0.9}{
            \makecell[{{p{\linewidth}}}]{
                \texttt{\tiny{[GM$|$GM]}}
                \texttt{RUCH: A, (4, 2)} \\
            }
        }
    }
    & & \\ \\

    \theutterance \stepcounter{utterance}  
    & & & \multicolumn{2}{p{0.3\linewidth}}{
        \cellcolor[rgb]{0.9,0.9,0.9}{
            \makecell[{{p{\linewidth}}}]{
                \texttt{\tiny{[GM$|$GM]}}
                \texttt{A moved to (4, 2)} \\
            }
        }
    }
    & & \\ \\

    \theutterance \stepcounter{utterance}  
    & & \multicolumn{4}{p{0.6\linewidth}}{
        \cellcolor[rgb]{0.9,0.9,0.9}{
            \makecell[{{p{\linewidth}}}]{
                \texttt{\tiny{[P1$\langle$GM]}}
                \texttt{Pomyślnie przeniesiono 'M' do (7,6). Twoja zaktualizowana siatka wygląda następująco:} \\
\texttt{```} \\
\texttt{ 1234567} \\
\texttt{╔════╤══╗} \\
\texttt{║◌◌◌O│S◌║ 1} \\
\texttt{║P◌◌◌├──╢ 2} \\
\texttt{║◌T◌A│◌◌║ 3} \\
\texttt{║◌┌──┼──╢ 4} \\
\texttt{║◌│◌◌│◌◌║ 5} \\
\texttt{╟─┴┬─┤◌M║ 6} \\
\texttt{║◌◌│◌│D◌║ 7} \\
\texttt{╚══╧═╧══╝} \\
\texttt{```} \\
\texttt{Obecnie grasz rundę 7 z maksymalnie 28.} \\
\texttt{Łącznie zgromadziliście 8 z 16 kar.} \\
\texttt{Drugi gracz przesunął obiekt na swojej siatce.} \\
\\ 
\\ 
\texttt{Jakie jest twoje następne polecenie?} \\
            }
        }
    }
    & \\ \\

    \theutterance \stepcounter{utterance}  
    & \multicolumn{4}{p{0.6\linewidth}}{
        \cellcolor[rgb]{0.9,0.9,0.9}{
            \makecell[{{p{\linewidth}}}]{
                \texttt{\tiny{[P1$\rangle$GM]}}
                \texttt{RUCH: D, (7, 5)} \\
            }
        }
    }
    & & \\ \\

    \theutterance \stepcounter{utterance}  
    & & & \multicolumn{2}{p{0.3\linewidth}}{
        \cellcolor[rgb]{0.9,0.9,0.9}{
            \makecell[{{p{\linewidth}}}]{
                \texttt{\tiny{[GM$|$GM]}}
                \texttt{RUCH: D, (7, 5)} \\
            }
        }
    }
    & & \\ \\

    \theutterance \stepcounter{utterance}  
    & & & \multicolumn{2}{p{0.3\linewidth}}{
        \cellcolor[rgb]{0.9,0.9,0.9}{
            \makecell[{{p{\linewidth}}}]{
                \texttt{\tiny{[GM$|$GM]}}
                \texttt{D moved to (7, 5)} \\
            }
        }
    }
    & & \\ \\

    \theutterance \stepcounter{utterance}  
    & & \multicolumn{4}{p{0.6\linewidth}}{
        \cellcolor[rgb]{0.9,0.9,0.9}{
            \makecell[{{p{\linewidth}}}]{
                \texttt{\tiny{[GM$\rangle$P2]}}
                \texttt{Pomyślnie przeniesiono 'A' do (4,2). Twoja zaktualizowana siatka wygląda następująco:} \\
\texttt{```} \\
\texttt{ 1234567} \\
\texttt{╔════╤══╗} \\
\texttt{║◌P◌◌│◌◌║ 1} \\
\texttt{║◌◌MA├──╢ 2} \\
\texttt{║◌◌S◌│◌◌║ 3} \\
\texttt{║O┌──┼──╢ 4} \\
\texttt{║◌│T◌│D◌║ 5} \\
\texttt{╟─┴┬─┤◌◌║ 6} \\
\texttt{║◌◌│◌│◌◌║ 7} \\
\texttt{╚══╧═╧══╝} \\
\texttt{```} \\
\texttt{Obecnie grasz rundę 8 z maksymalnie 28.} \\
\texttt{Łącznie zgromadziliście 8 z 16 kar.} \\
\texttt{Drugi gracz przesunął obiekt na swojej siatce.} \\
\\ 
\\ 
\texttt{Jakie jest twoje następne polecenie?} \\
            }
        }
    }
    & \\ \\

    \theutterance \stepcounter{utterance}  
    & & & \multicolumn{4}{p{0.6\linewidth}}{
        \cellcolor[rgb]{0.9,0.9,0.9}{
            \makecell[{{p{\linewidth}}}]{
                \texttt{\tiny{[GM$\langle$P2]}}
                \texttt{RUCH: M, (2, 3)} \\
            }
        }
    }
     \\ \\

    \theutterance \stepcounter{utterance}  
    & & & \multicolumn{2}{p{0.3\linewidth}}{
        \cellcolor[rgb]{0.9,0.9,0.9}{
            \makecell[{{p{\linewidth}}}]{
                \texttt{\tiny{[GM$|$GM]}}
                \texttt{RUCH: M, (2, 3)} \\
            }
        }
    }
    & & \\ \\

    \theutterance \stepcounter{utterance}  
    & & & \multicolumn{2}{p{0.3\linewidth}}{
        \cellcolor[rgb]{0.9,0.9,0.9}{
            \makecell[{{p{\linewidth}}}]{
                \texttt{\tiny{[GM$|$GM]}}
                \texttt{M moved to (2, 3)} \\
            }
        }
    }
    & & \\ \\

    \theutterance \stepcounter{utterance}  
    & & \multicolumn{4}{p{0.6\linewidth}}{
        \cellcolor[rgb]{0.9,0.9,0.9}{
            \makecell[{{p{\linewidth}}}]{
                \texttt{\tiny{[P1$\langle$GM]}}
                \texttt{Pomyślnie przeniesiono 'D' do (7,5). Twoja zaktualizowana siatka wygląda następująco:} \\
\texttt{```} \\
\texttt{ 1234567} \\
\texttt{╔════╤══╗} \\
\texttt{║◌◌◌O│S◌║ 1} \\
\texttt{║P◌◌◌├──╢ 2} \\
\texttt{║◌T◌A│◌◌║ 3} \\
\texttt{║◌┌──┼──╢ 4} \\
\texttt{║◌│◌◌│◌D║ 5} \\
\texttt{╟─┴┬─┤◌M║ 6} \\
\texttt{║◌◌│◌│◌◌║ 7} \\
\texttt{╚══╧═╧══╝} \\
\texttt{```} \\
\texttt{Obecnie grasz rundę 8 z maksymalnie 28.} \\
\texttt{Łącznie zgromadziliście 8 z 16 kar.} \\
\texttt{Drugi gracz przesunął obiekt na swojej siatce.} \\
\\ 
\\ 
\texttt{Jakie jest twoje następne polecenie?} \\
            }
        }
    }
    & \\ \\

    \theutterance \stepcounter{utterance}  
    & \multicolumn{4}{p{0.6\linewidth}}{
        \cellcolor[rgb]{0.9,0.9,0.9}{
            \makecell[{{p{\linewidth}}}]{
                \texttt{\tiny{[P1$\rangle$GM]}}
                \texttt{POWIEDZ: Potrzebuję pomocy w ustaleniu finalnej pozycji każdego obiektu na obu planszach.} \\
            }
        }
    }
    & & \\ \\

    \theutterance \stepcounter{utterance}  
    & & & \multicolumn{2}{p{0.3\linewidth}}{
        \cellcolor[rgb]{0.9,0.9,0.9}{
            \makecell[{{p{\linewidth}}}]{
                \texttt{\tiny{[GM$|$GM]}}
                \texttt{POWIEDZ: Potrzebuję pomocy w ustaleniu finalnej pozycji każdego obiektu na obu planszach.} \\
            }
        }
    }
    & & \\ \\

    \theutterance \stepcounter{utterance}  
    & & \multicolumn{4}{p{0.6\linewidth}}{
        \cellcolor[rgb]{0.9,0.9,0.9}{
            \makecell[{{p{\linewidth}}}]{
                \texttt{\tiny{[GM$\rangle$P2]}}
                \texttt{Pomyślnie przeniesiono 'M' do (2,3). Twoja zaktualizowana siatka wygląda następująco:} \\
\texttt{```} \\
\texttt{ 1234567} \\
\texttt{╔════╤══╗} \\
\texttt{║◌P◌◌│◌◌║ 1} \\
\texttt{║◌◌◌A├──╢ 2} \\
\texttt{║◌MS◌│◌◌║ 3} \\
\texttt{║O┌──┼──╢ 4} \\
\texttt{║◌│T◌│D◌║ 5} \\
\texttt{╟─┴┬─┤◌◌║ 6} \\
\texttt{║◌◌│◌│◌◌║ 7} \\
\texttt{╚══╧═╧══╝} \\
\texttt{```} \\
\texttt{Obecnie grasz rundę 9 z maksymalnie 28.} \\
\texttt{Łącznie zgromadziliście 8 z 16 kar.} \\
\texttt{Drugi gracz wysłał tę wiadomość:} \\
\\ 
\\ 
\texttt{"Potrzebuję pomocy w ustaleniu finalnej pozycji każdego obiektu na obu planszach."} \\
\\ 
\\ 
\texttt{Jakie jest twoje następne polecenie?} \\
            }
        }
    }
    & \\ \\

    \theutterance \stepcounter{utterance}  
    & & & \multicolumn{4}{p{0.6\linewidth}}{
        \cellcolor[rgb]{0.9,0.9,0.9}{
            \makecell[{{p{\linewidth}}}]{
                \texttt{\tiny{[GM$\langle$P2]}}
                \texttt{POWIEDZ: Oto nasza bieżąca konfiguracja docelowa: O(1,4), T(3,5), S(3,3), D(6,5), A(4,2), M(2,3). Poproś o potwierdzenie lub ustalenie innych pozycji.} \\
            }
        }
    }
     \\ \\

    \theutterance \stepcounter{utterance}  
    & & & \multicolumn{2}{p{0.3\linewidth}}{
        \cellcolor[rgb]{0.9,0.9,0.9}{
            \makecell[{{p{\linewidth}}}]{
                \texttt{\tiny{[GM$|$GM]}}
                \texttt{POWIEDZ: Oto nasza bieżąca konfiguracja docelowa: O(1,4), T(3,5), S(3,3), D(6,5), A(4,2), M(2,3). Poproś o potwierdzenie lub ustalenie innych pozycji.} \\
            }
        }
    }
    & & \\ \\

    \theutterance \stepcounter{utterance}  
    & & \multicolumn{4}{p{0.6\linewidth}}{
        \cellcolor[rgb]{0.9,0.9,0.9}{
            \makecell[{{p{\linewidth}}}]{
                \texttt{\tiny{[P1$\langle$GM]}}
                \texttt{Twoja wiadomość została przekazana drugiemu graczowi.} \\
\\ 
\\ 
\texttt{Obecnie grasz rundę 9 z maksymalnie 28.} \\
\texttt{Łącznie zgromadziliście 8 z 16 kar.} \\
\texttt{Drugi gracz wysłał tę wiadomość:} \\
\\ 
\\ 
\texttt{"Oto nasza bieżąca konfiguracja docelowa: O(1,4), T(3,5), S(3,3), D(6,5), A(4,2), M(2,3). Poproś o potwierdzenie lub ustalenie innych pozycji."} \\
\\ 
\\ 
\texttt{Jakie jest twoje następne polecenie?} \\
            }
        }
    }
    & \\ \\

    \theutterance \stepcounter{utterance}  
    & \multicolumn{4}{p{0.6\linewidth}}{
        \cellcolor[rgb]{0.9,0.9,0.9}{
            \makecell[{{p{\linewidth}}}]{
                \texttt{\tiny{[P1$\rangle$GM]}}
                \texttt{POWIEDZ: Potwierdzam uzgodnioną konfigurację docelową: O(1,4), T(3,5), S(3,3), D(6,5), A(4,2), M(2,3).} \\
            }
        }
    }
    & & \\ \\

    \theutterance \stepcounter{utterance}  
    & & & \multicolumn{2}{p{0.3\linewidth}}{
        \cellcolor[rgb]{0.9,0.9,0.9}{
            \makecell[{{p{\linewidth}}}]{
                \texttt{\tiny{[GM$|$GM]}}
                \texttt{POWIEDZ: Potwierdzam uzgodnioną konfigurację docelową: O(1,4), T(3,5), S(3,3), D(6,5), A(4,2), M(2,3).} \\
            }
        }
    }
    & & \\ \\

    \theutterance \stepcounter{utterance}  
    & & \multicolumn{4}{p{0.6\linewidth}}{
        \cellcolor[rgb]{0.9,0.9,0.9}{
            \makecell[{{p{\linewidth}}}]{
                \texttt{\tiny{[GM$\rangle$P2]}}
                \texttt{Twoja wiadomość została przekazana drugiemu graczowi.} \\
\\ 
\\ 
\texttt{Obecnie grasz rundę 10 z maksymalnie 28.} \\
\texttt{Łącznie zgromadziliście 8 z 16 kar.} \\
\texttt{Drugi gracz wysłał tę wiadomość:} \\
\\ 
\\ 
\texttt{"Potwierdzam uzgodnioną konfigurację docelową: O(1,4), T(3,5), S(3,3), D(6,5), A(4,2), M(2,3)."} \\
\\ 
\\ 
\texttt{Jakie jest twoje następne polecenie?} \\
            }
        }
    }
    & \\ \\

    \theutterance \stepcounter{utterance}  
    & & & \multicolumn{4}{p{0.6\linewidth}}{
        \cellcolor[rgb]{0.9,0.9,0.9}{
            \makecell[{{p{\linewidth}}}]{
                \texttt{\tiny{[GM$\langle$P2]}}
                \texttt{RUCH: A, (4, 2)} \\
            }
        }
    }
     \\ \\

    \theutterance \stepcounter{utterance}  
    & & & \multicolumn{2}{p{0.3\linewidth}}{
        \cellcolor[rgb]{0.9,0.9,0.9}{
            \makecell[{{p{\linewidth}}}]{
                \texttt{\tiny{[GM$|$GM]}}
                \texttt{RUCH: A, (4, 2)} \\
            }
        }
    }
    & & \\ \\

    \theutterance \stepcounter{utterance}  
    & & & \multicolumn{2}{p{0.3\linewidth}}{
        \cellcolor[rgb]{0.9,0.9,0.9}{
            \makecell[{{p{\linewidth}}}]{
                \texttt{\tiny{[GM$|$GM]}}
                \texttt{Kara: (4,2) nie jest puste, ale zawiera 'A'.} \\
\\ 
\texttt{Łącznie zgromadziliście 9 z 16 kar.} \\
\texttt{Spróbuj ponownie!} \\
            }
        }
    }
    & & \\ \\

    \theutterance \stepcounter{utterance}  
    & & \multicolumn{4}{p{0.6\linewidth}}{
        \cellcolor[rgb]{0.9,0.9,0.9}{
            \makecell[{{p{\linewidth}}}]{
                \texttt{\tiny{[GM$\rangle$P2]}}
                \texttt{Kara: (4,2) nie jest puste, ale zawiera 'A'.} \\
\\ 
\texttt{Łącznie zgromadziliście 9 z 16 kar.} \\
\texttt{Spróbuj ponownie!} \\
            }
        }
    }
    & \\ \\

    \theutterance \stepcounter{utterance}  
    & & & \multicolumn{4}{p{0.6\linewidth}}{
        \cellcolor[rgb]{0.9,0.9,0.9}{
            \makecell[{{p{\linewidth}}}]{
                \texttt{\tiny{[GM$\langle$P2]}}
                \texttt{RUCH: A, (3, 2)} \\
            }
        }
    }
     \\ \\

    \theutterance \stepcounter{utterance}  
    & & & \multicolumn{2}{p{0.3\linewidth}}{
        \cellcolor[rgb]{0.9,0.9,0.9}{
            \makecell[{{p{\linewidth}}}]{
                \texttt{\tiny{[GM$|$GM]}}
                \texttt{RUCH: A, (3, 2)} \\
            }
        }
    }
    & & \\ \\

    \theutterance \stepcounter{utterance}  
    & & & \multicolumn{2}{p{0.3\linewidth}}{
        \cellcolor[rgb]{0.9,0.9,0.9}{
            \makecell[{{p{\linewidth}}}]{
                \texttt{\tiny{[GM$|$GM]}}
                \texttt{A moved to (3, 2)} \\
            }
        }
    }
    & & \\ \\

    \theutterance \stepcounter{utterance}  
    & & \multicolumn{4}{p{0.6\linewidth}}{
        \cellcolor[rgb]{0.9,0.9,0.9}{
            \makecell[{{p{\linewidth}}}]{
                \texttt{\tiny{[P1$\langle$GM]}}
                \texttt{Twoja wiadomość została przekazana drugiemu graczowi.} \\
\\ 
\\ 
\texttt{Obecnie grasz rundę 10 z maksymalnie 28.} \\
\texttt{Łącznie zgromadziliście 9 z 16 kar.} \\
\texttt{Drugi gracz przesunął obiekt na swojej siatce.} \\
\\ 
\\ 
\texttt{Jakie jest twoje następne polecenie?} \\
            }
        }
    }
    & \\ \\

    \theutterance \stepcounter{utterance}  
    & \multicolumn{4}{p{0.6\linewidth}}{
        \cellcolor[rgb]{0.9,0.9,0.9}{
            \makecell[{{p{\linewidth}}}]{
                \texttt{\tiny{[P1$\rangle$GM]}}
                \texttt{RUCH: O, (1, 4)} \\
            }
        }
    }
    & & \\ \\

    \theutterance \stepcounter{utterance}  
    & & & \multicolumn{2}{p{0.3\linewidth}}{
        \cellcolor[rgb]{0.9,0.9,0.9}{
            \makecell[{{p{\linewidth}}}]{
                \texttt{\tiny{[GM$|$GM]}}
                \texttt{RUCH: O, (1, 4)} \\
            }
        }
    }
    & & \\ \\

    \theutterance \stepcounter{utterance}  
    & & & \multicolumn{2}{p{0.3\linewidth}}{
        \cellcolor[rgb]{0.9,0.9,0.9}{
            \makecell[{{p{\linewidth}}}]{
                \texttt{\tiny{[GM$|$GM]}}
                \texttt{O moved to (1, 4)} \\
            }
        }
    }
    & & \\ \\

    \theutterance \stepcounter{utterance}  
    & & \multicolumn{4}{p{0.6\linewidth}}{
        \cellcolor[rgb]{0.9,0.9,0.9}{
            \makecell[{{p{\linewidth}}}]{
                \texttt{\tiny{[GM$\rangle$P2]}}
                \texttt{Pomyślnie przeniesiono 'A' do (3,2). Twoja zaktualizowana siatka wygląda następująco:} \\
\texttt{```} \\
\texttt{ 1234567} \\
\texttt{╔════╤══╗} \\
\texttt{║◌P◌◌│◌◌║ 1} \\
\texttt{║◌◌A◌├──╢ 2} \\
\texttt{║◌MS◌│◌◌║ 3} \\
\texttt{║O┌──┼──╢ 4} \\
\texttt{║◌│T◌│D◌║ 5} \\
\texttt{╟─┴┬─┤◌◌║ 6} \\
\texttt{║◌◌│◌│◌◌║ 7} \\
\texttt{╚══╧═╧══╝} \\
\texttt{```} \\
\texttt{Obecnie grasz rundę 11 z maksymalnie 28.} \\
\texttt{Łącznie zgromadziliście 9 z 16 kar.} \\
\texttt{Drugi gracz przesunął obiekt na swojej siatce.} \\
\\ 
\\ 
\texttt{Jakie jest twoje następne polecenie?} \\
            }
        }
    }
    & \\ \\

    \theutterance \stepcounter{utterance}  
    & & & \multicolumn{4}{p{0.6\linewidth}}{
        \cellcolor[rgb]{0.9,0.9,0.9}{
            \makecell[{{p{\linewidth}}}]{
                \texttt{\tiny{[GM$\langle$P2]}}
                \texttt{RUCH: M, (2, 3)} \\
            }
        }
    }
     \\ \\

    \theutterance \stepcounter{utterance}  
    & & & \multicolumn{2}{p{0.3\linewidth}}{
        \cellcolor[rgb]{0.9,0.9,0.9}{
            \makecell[{{p{\linewidth}}}]{
                \texttt{\tiny{[GM$|$GM]}}
                \texttt{RUCH: M, (2, 3)} \\
            }
        }
    }
    & & \\ \\

    \theutterance \stepcounter{utterance}  
    & & & \multicolumn{2}{p{0.3\linewidth}}{
        \cellcolor[rgb]{0.9,0.9,0.9}{
            \makecell[{{p{\linewidth}}}]{
                \texttt{\tiny{[GM$|$GM]}}
                \texttt{Kara: (2,3) nie jest puste, ale zawiera 'M'.} \\
\\ 
\texttt{Łącznie zgromadziliście 10 z 16 kar.} \\
\texttt{Spróbuj ponownie!} \\
            }
        }
    }
    & & \\ \\

    \theutterance \stepcounter{utterance}  
    & & \multicolumn{4}{p{0.6\linewidth}}{
        \cellcolor[rgb]{0.9,0.9,0.9}{
            \makecell[{{p{\linewidth}}}]{
                \texttt{\tiny{[GM$\rangle$P2]}}
                \texttt{Kara: (2,3) nie jest puste, ale zawiera 'M'.} \\
\\ 
\texttt{Łącznie zgromadziliście 10 z 16 kar.} \\
\texttt{Spróbuj ponownie!} \\
            }
        }
    }
    & \\ \\

    \theutterance \stepcounter{utterance}  
    & & & \multicolumn{4}{p{0.6\linewidth}}{
        \cellcolor[rgb]{0.9,0.9,0.9}{
            \makecell[{{p{\linewidth}}}]{
                \texttt{\tiny{[GM$\langle$P2]}}
                \texttt{RUCH: M, (2, 4)} \\
            }
        }
    }
     \\ \\

    \theutterance \stepcounter{utterance}  
    & & & \multicolumn{2}{p{0.3\linewidth}}{
        \cellcolor[rgb]{0.9,0.9,0.9}{
            \makecell[{{p{\linewidth}}}]{
                \texttt{\tiny{[GM$|$GM]}}
                \texttt{RUCH: M, (2, 4)} \\
            }
        }
    }
    & & \\ \\

    \theutterance \stepcounter{utterance}  
    & & & \multicolumn{2}{p{0.3\linewidth}}{
        \cellcolor[rgb]{0.9,0.9,0.9}{
            \makecell[{{p{\linewidth}}}]{
                \texttt{\tiny{[GM$|$GM]}}
                \texttt{Kara: (2,4) nie jest puste, ale zawiera '┌'.} \\
\\ 
\texttt{Łącznie zgromadziliście 11 z 16 kar.} \\
\texttt{Spróbuj ponownie!} \\
            }
        }
    }
    & & \\ \\

    \theutterance \stepcounter{utterance}  
    & & \multicolumn{4}{p{0.6\linewidth}}{
        \cellcolor[rgb]{0.9,0.9,0.9}{
            \makecell[{{p{\linewidth}}}]{
                \texttt{\tiny{[GM$\rangle$P2]}}
                \texttt{Kara: (2,4) nie jest puste, ale zawiera '┌'.} \\
\\ 
\texttt{Łącznie zgromadziliście 11 z 16 kar.} \\
\texttt{Spróbuj ponownie!} \\
            }
        }
    }
    & \\ \\

    \theutterance \stepcounter{utterance}  
    & & & \multicolumn{4}{p{0.6\linewidth}}{
        \cellcolor[rgb]{0.9,0.9,0.9}{
            \makecell[{{p{\linewidth}}}]{
                \texttt{\tiny{[GM$\langle$P2]}}
                \texttt{RUCH: M, (2, 5)} \\
            }
        }
    }
     \\ \\

    \theutterance \stepcounter{utterance}  
    & & & \multicolumn{2}{p{0.3\linewidth}}{
        \cellcolor[rgb]{0.9,0.9,0.9}{
            \makecell[{{p{\linewidth}}}]{
                \texttt{\tiny{[GM$|$GM]}}
                \texttt{RUCH: M, (2, 5)} \\
            }
        }
    }
    & & \\ \\

    \theutterance \stepcounter{utterance}  
    & & & \multicolumn{2}{p{0.3\linewidth}}{
        \cellcolor[rgb]{0.9,0.9,0.9}{
            \makecell[{{p{\linewidth}}}]{
                \texttt{\tiny{[GM$|$GM]}}
                \texttt{Kara: (2,5) nie jest puste, ale zawiera '│'.} \\
\\ 
\texttt{Łącznie zgromadziliście 12 z 16 kar.} \\
\texttt{Spróbuj ponownie!} \\
            }
        }
    }
    & & \\ \\

    \theutterance \stepcounter{utterance}  
    & & \multicolumn{4}{p{0.6\linewidth}}{
        \cellcolor[rgb]{0.9,0.9,0.9}{
            \makecell[{{p{\linewidth}}}]{
                \texttt{\tiny{[GM$\rangle$P2]}}
                \texttt{Kara: (2,5) nie jest puste, ale zawiera '│'.} \\
\\ 
\texttt{Łącznie zgromadziliście 12 z 16 kar.} \\
\texttt{Spróbuj ponownie!} \\
            }
        }
    }
    & \\ \\

    \theutterance \stepcounter{utterance}  
    & & & \multicolumn{4}{p{0.6\linewidth}}{
        \cellcolor[rgb]{0.9,0.9,0.9}{
            \makecell[{{p{\linewidth}}}]{
                \texttt{\tiny{[GM$\langle$P2]}}
                \texttt{RUCH: M, (1, 3)} \\
            }
        }
    }
     \\ \\

    \theutterance \stepcounter{utterance}  
    & & & \multicolumn{2}{p{0.3\linewidth}}{
        \cellcolor[rgb]{0.9,0.9,0.9}{
            \makecell[{{p{\linewidth}}}]{
                \texttt{\tiny{[GM$|$GM]}}
                \texttt{RUCH: M, (1, 3)} \\
            }
        }
    }
    & & \\ \\

    \theutterance \stepcounter{utterance}  
    & & & \multicolumn{2}{p{0.3\linewidth}}{
        \cellcolor[rgb]{0.9,0.9,0.9}{
            \makecell[{{p{\linewidth}}}]{
                \texttt{\tiny{[GM$|$GM]}}
                \texttt{M moved to (1, 3)} \\
            }
        }
    }
    & & \\ \\

    \theutterance \stepcounter{utterance}  
    & & \multicolumn{4}{p{0.6\linewidth}}{
        \cellcolor[rgb]{0.9,0.9,0.9}{
            \makecell[{{p{\linewidth}}}]{
                \texttt{\tiny{[P1$\langle$GM]}}
                \texttt{Pomyślnie przeniesiono 'O' do (1,4). Twoja zaktualizowana siatka wygląda następująco:} \\
\texttt{```} \\
\texttt{ 1234567} \\
\texttt{╔════╤══╗} \\
\texttt{║◌◌◌◌│S◌║ 1} \\
\texttt{║P◌◌◌├──╢ 2} \\
\texttt{║◌T◌A│◌◌║ 3} \\
\texttt{║O┌──┼──╢ 4} \\
\texttt{║◌│◌◌│◌D║ 5} \\
\texttt{╟─┴┬─┤◌M║ 6} \\
\texttt{║◌◌│◌│◌◌║ 7} \\
\texttt{╚══╧═╧══╝} \\
\texttt{```} \\
\texttt{Obecnie grasz rundę 11 z maksymalnie 28.} \\
\texttt{Łącznie zgromadziliście 12 z 16 kar.} \\
\texttt{Drugi gracz przesunął obiekt na swojej siatce.} \\
\\ 
\\ 
\texttt{Jakie jest twoje następne polecenie?} \\
            }
        }
    }
    & \\ \\

    \theutterance \stepcounter{utterance}  
    & \multicolumn{4}{p{0.6\linewidth}}{
        \cellcolor[rgb]{0.9,0.9,0.9}{
            \makecell[{{p{\linewidth}}}]{
                \texttt{\tiny{[P1$\rangle$GM]}}
                \texttt{RUCH: T, (3, 5)} \\
            }
        }
    }
    & & \\ \\

    \theutterance \stepcounter{utterance}  
    & & & \multicolumn{2}{p{0.3\linewidth}}{
        \cellcolor[rgb]{0.9,0.9,0.9}{
            \makecell[{{p{\linewidth}}}]{
                \texttt{\tiny{[GM$|$GM]}}
                \texttt{RUCH: T, (3, 5)} \\
            }
        }
    }
    & & \\ \\

    \theutterance \stepcounter{utterance}  
    & & & \multicolumn{2}{p{0.3\linewidth}}{
        \cellcolor[rgb]{0.9,0.9,0.9}{
            \makecell[{{p{\linewidth}}}]{
                \texttt{\tiny{[GM$|$GM]}}
                \texttt{T moved to (3, 5)} \\
            }
        }
    }
    & & \\ \\

    \theutterance \stepcounter{utterance}  
    & & \multicolumn{4}{p{0.6\linewidth}}{
        \cellcolor[rgb]{0.9,0.9,0.9}{
            \makecell[{{p{\linewidth}}}]{
                \texttt{\tiny{[GM$\rangle$P2]}}
                \texttt{Pomyślnie przeniesiono 'M' do (1,3). Twoja zaktualizowana siatka wygląda następująco:} \\
\texttt{```} \\
\texttt{ 1234567} \\
\texttt{╔════╤══╗} \\
\texttt{║◌P◌◌│◌◌║ 1} \\
\texttt{║◌◌A◌├──╢ 2} \\
\texttt{║M◌S◌│◌◌║ 3} \\
\texttt{║O┌──┼──╢ 4} \\
\texttt{║◌│T◌│D◌║ 5} \\
\texttt{╟─┴┬─┤◌◌║ 6} \\
\texttt{║◌◌│◌│◌◌║ 7} \\
\texttt{╚══╧═╧══╝} \\
\texttt{```} \\
\texttt{Obecnie grasz rundę 12 z maksymalnie 28.} \\
\texttt{Łącznie zgromadziliście 12 z 16 kar.} \\
\texttt{Drugi gracz przesunął obiekt na swojej siatce.} \\
\\ 
\\ 
\texttt{Jakie jest twoje następne polecenie?} \\
            }
        }
    }
    & \\ \\

    \theutterance \stepcounter{utterance}  
    & & & \multicolumn{4}{p{0.6\linewidth}}{
        \cellcolor[rgb]{0.9,0.9,0.9}{
            \makecell[{{p{\linewidth}}}]{
                \texttt{\tiny{[GM$\langle$P2]}}
                \texttt{RUCH: S, (3, 2)} \\
            }
        }
    }
     \\ \\

    \theutterance \stepcounter{utterance}  
    & & & \multicolumn{2}{p{0.3\linewidth}}{
        \cellcolor[rgb]{0.9,0.9,0.9}{
            \makecell[{{p{\linewidth}}}]{
                \texttt{\tiny{[GM$|$GM]}}
                \texttt{RUCH: S, (3, 2)} \\
            }
        }
    }
    & & \\ \\

    \theutterance \stepcounter{utterance}  
    & & & \multicolumn{2}{p{0.3\linewidth}}{
        \cellcolor[rgb]{0.9,0.9,0.9}{
            \makecell[{{p{\linewidth}}}]{
                \texttt{\tiny{[GM$|$GM]}}
                \texttt{Kara: (3,2) nie jest puste, ale zawiera 'A'.} \\
\\ 
\texttt{Łącznie zgromadziliście 13 z 16 kar.} \\
\texttt{Spróbuj ponownie!} \\
            }
        }
    }
    & & \\ \\

    \theutterance \stepcounter{utterance}  
    & & \multicolumn{4}{p{0.6\linewidth}}{
        \cellcolor[rgb]{0.9,0.9,0.9}{
            \makecell[{{p{\linewidth}}}]{
                \texttt{\tiny{[GM$\rangle$P2]}}
                \texttt{Kara: (3,2) nie jest puste, ale zawiera 'A'.} \\
\\ 
\texttt{Łącznie zgromadziliście 13 z 16 kar.} \\
\texttt{Spróbuj ponownie!} \\
            }
        }
    }
    & \\ \\

    \theutterance \stepcounter{utterance}  
    & & & \multicolumn{4}{p{0.6\linewidth}}{
        \cellcolor[rgb]{0.9,0.9,0.9}{
            \makecell[{{p{\linewidth}}}]{
                \texttt{\tiny{[GM$\langle$P2]}}
                \texttt{RUCH: S, (3, 3)} \\
            }
        }
    }
     \\ \\

    \theutterance \stepcounter{utterance}  
    & & & \multicolumn{2}{p{0.3\linewidth}}{
        \cellcolor[rgb]{0.9,0.9,0.9}{
            \makecell[{{p{\linewidth}}}]{
                \texttt{\tiny{[GM$|$GM]}}
                \texttt{RUCH: S, (3, 3)} \\
            }
        }
    }
    & & \\ \\

    \theutterance \stepcounter{utterance}  
    & & & \multicolumn{2}{p{0.3\linewidth}}{
        \cellcolor[rgb]{0.9,0.9,0.9}{
            \makecell[{{p{\linewidth}}}]{
                \texttt{\tiny{[GM$|$GM]}}
                \texttt{Kara: (3,3) nie jest puste, ale zawiera 'S'.} \\
\\ 
\texttt{Łącznie zgromadziliście 14 z 16 kar.} \\
\texttt{Spróbuj ponownie!} \\
            }
        }
    }
    & & \\ \\

    \theutterance \stepcounter{utterance}  
    & & \multicolumn{4}{p{0.6\linewidth}}{
        \cellcolor[rgb]{0.9,0.9,0.9}{
            \makecell[{{p{\linewidth}}}]{
                \texttt{\tiny{[GM$\rangle$P2]}}
                \texttt{Kara: (3,3) nie jest puste, ale zawiera 'S'.} \\
\\ 
\texttt{Łącznie zgromadziliście 14 z 16 kar.} \\
\texttt{Spróbuj ponownie!} \\
            }
        }
    }
    & \\ \\

    \theutterance \stepcounter{utterance}  
    & & & \multicolumn{4}{p{0.6\linewidth}}{
        \cellcolor[rgb]{0.9,0.9,0.9}{
            \makecell[{{p{\linewidth}}}]{
                \texttt{\tiny{[GM$\langle$P2]}}
                \texttt{RUCH: S, (4, 2)} \\
            }
        }
    }
     \\ \\

    \theutterance \stepcounter{utterance}  
    & & & \multicolumn{2}{p{0.3\linewidth}}{
        \cellcolor[rgb]{0.9,0.9,0.9}{
            \makecell[{{p{\linewidth}}}]{
                \texttt{\tiny{[GM$|$GM]}}
                \texttt{RUCH: S, (4, 2)} \\
            }
        }
    }
    & & \\ \\

    \theutterance \stepcounter{utterance}  
    & & & \multicolumn{2}{p{0.3\linewidth}}{
        \cellcolor[rgb]{0.9,0.9,0.9}{
            \makecell[{{p{\linewidth}}}]{
                \texttt{\tiny{[GM$|$GM]}}
                \texttt{S moved to (4, 2)} \\
            }
        }
    }
    & & \\ \\

    \theutterance \stepcounter{utterance}  
    & & \multicolumn{4}{p{0.6\linewidth}}{
        \cellcolor[rgb]{0.9,0.9,0.9}{
            \makecell[{{p{\linewidth}}}]{
                \texttt{\tiny{[P1$\langle$GM]}}
                \texttt{Pomyślnie przeniesiono 'T' do (3,5). Twoja zaktualizowana siatka wygląda następująco:} \\
\texttt{```} \\
\texttt{ 1234567} \\
\texttt{╔════╤══╗} \\
\texttt{║◌◌◌◌│S◌║ 1} \\
\texttt{║P◌◌◌├──╢ 2} \\
\texttt{║◌◌◌A│◌◌║ 3} \\
\texttt{║O┌──┼──╢ 4} \\
\texttt{║◌│T◌│◌D║ 5} \\
\texttt{╟─┴┬─┤◌M║ 6} \\
\texttt{║◌◌│◌│◌◌║ 7} \\
\texttt{╚══╧═╧══╝} \\
\texttt{```} \\
\texttt{Obecnie grasz rundę 12 z maksymalnie 28.} \\
\texttt{Łącznie zgromadziliście 14 z 16 kar.} \\
\texttt{Drugi gracz przesunął obiekt na swojej siatce.} \\
\\ 
\\ 
\texttt{Jakie jest twoje następne polecenie?} \\
            }
        }
    }
    & \\ \\

    \theutterance \stepcounter{utterance}  
    & \multicolumn{4}{p{0.6\linewidth}}{
        \cellcolor[rgb]{0.9,0.9,0.9}{
            \makecell[{{p{\linewidth}}}]{
                \texttt{\tiny{[P1$\rangle$GM]}}
                \texttt{RUCH: S, (3, 3)} \\
            }
        }
    }
    & & \\ \\

    \theutterance \stepcounter{utterance}  
    & & & \multicolumn{2}{p{0.3\linewidth}}{
        \cellcolor[rgb]{0.9,0.9,0.9}{
            \makecell[{{p{\linewidth}}}]{
                \texttt{\tiny{[GM$|$GM]}}
                \texttt{RUCH: S, (3, 3)} \\
            }
        }
    }
    & & \\ \\

    \theutterance \stepcounter{utterance}  
    & & & \multicolumn{2}{p{0.3\linewidth}}{
        \cellcolor[rgb]{0.9,0.9,0.9}{
            \makecell[{{p{\linewidth}}}]{
                \texttt{\tiny{[GM$|$GM]}}
                \texttt{S moved to (3, 3)} \\
            }
        }
    }
    & & \\ \\

    \theutterance \stepcounter{utterance}  
    & & \multicolumn{4}{p{0.6\linewidth}}{
        \cellcolor[rgb]{0.9,0.9,0.9}{
            \makecell[{{p{\linewidth}}}]{
                \texttt{\tiny{[GM$\rangle$P2]}}
                \texttt{Pomyślnie przeniesiono 'S' do (4,2). Twoja zaktualizowana siatka wygląda następująco:} \\
\texttt{```} \\
\texttt{ 1234567} \\
\texttt{╔════╤══╗} \\
\texttt{║◌P◌◌│◌◌║ 1} \\
\texttt{║◌◌AS├──╢ 2} \\
\texttt{║M◌◌◌│◌◌║ 3} \\
\texttt{║O┌──┼──╢ 4} \\
\texttt{║◌│T◌│D◌║ 5} \\
\texttt{╟─┴┬─┤◌◌║ 6} \\
\texttt{║◌◌│◌│◌◌║ 7} \\
\texttt{╚══╧═╧══╝} \\
\texttt{```} \\
\texttt{Obecnie grasz rundę 13 z maksymalnie 28.} \\
\texttt{Łącznie zgromadziliście 14 z 16 kar.} \\
\texttt{Drugi gracz przesunął obiekt na swojej siatce.} \\
\\ 
\\ 
\texttt{Jakie jest twoje następne polecenie?} \\
            }
        }
    }
    & \\ \\

    \theutterance \stepcounter{utterance}  
    & & & \multicolumn{4}{p{0.6\linewidth}}{
        \cellcolor[rgb]{0.9,0.9,0.9}{
            \makecell[{{p{\linewidth}}}]{
                \texttt{\tiny{[GM$\langle$P2]}}
                \texttt{RUCH: T, (5, 4)} \\
            }
        }
    }
     \\ \\

    \theutterance \stepcounter{utterance}  
    & & & \multicolumn{2}{p{0.3\linewidth}}{
        \cellcolor[rgb]{0.9,0.9,0.9}{
            \makecell[{{p{\linewidth}}}]{
                \texttt{\tiny{[GM$|$GM]}}
                \texttt{RUCH: T, (5, 4)} \\
            }
        }
    }
    & & \\ \\

    \theutterance \stepcounter{utterance}  
    & & & \multicolumn{2}{p{0.3\linewidth}}{
        \cellcolor[rgb]{0.9,0.9,0.9}{
            \makecell[{{p{\linewidth}}}]{
                \texttt{\tiny{[GM$|$GM]}}
                \texttt{Kara: (5,4) nie jest puste, ale zawiera '┼'.} \\
\\ 
\texttt{Łącznie zgromadziliście 15 z 16 kar.} \\
\texttt{Spróbuj ponownie!} \\
            }
        }
    }
    & & \\ \\

    \theutterance \stepcounter{utterance}  
    & & \multicolumn{4}{p{0.6\linewidth}}{
        \cellcolor[rgb]{0.9,0.9,0.9}{
            \makecell[{{p{\linewidth}}}]{
                \texttt{\tiny{[GM$\rangle$P2]}}
                \texttt{Kara: (5,4) nie jest puste, ale zawiera '┼'.} \\
\\ 
\texttt{Łącznie zgromadziliście 15 z 16 kar.} \\
\texttt{Spróbuj ponownie!} \\
            }
        }
    }
    & \\ \\

    \theutterance \stepcounter{utterance}  
    & & & \multicolumn{4}{p{0.6\linewidth}}{
        \cellcolor[rgb]{0.9,0.9,0.9}{
            \makecell[{{p{\linewidth}}}]{
                \texttt{\tiny{[GM$\langle$P2]}}
                \texttt{RUCH: T, (5, 5)} \\
            }
        }
    }
     \\ \\

    \theutterance \stepcounter{utterance}  
    & & & \multicolumn{2}{p{0.3\linewidth}}{
        \cellcolor[rgb]{0.9,0.9,0.9}{
            \makecell[{{p{\linewidth}}}]{
                \texttt{\tiny{[GM$|$GM]}}
                \texttt{RUCH: T, (5, 5)} \\
            }
        }
    }
    & & \\ \\

    \theutterance \stepcounter{utterance}  
    & & & \multicolumn{2}{p{0.3\linewidth}}{
        \cellcolor[rgb]{0.9,0.9,0.9}{
            \makecell[{{p{\linewidth}}}]{
                \texttt{\tiny{[GM$|$GM]}}
                \texttt{Kara: (5,5) nie jest puste, ale zawiera '│'.} \\
\\ 
\texttt{Łącznie zgromadziliście 16 z 16 kar.} \\
\texttt{Spróbuj ponownie!} \\
            }
        }
    }
    & & \\ \\

    \theutterance \stepcounter{utterance}  
    & & \multicolumn{4}{p{0.6\linewidth}}{
        \cellcolor[rgb]{0.9,0.9,0.9}{
            \makecell[{{p{\linewidth}}}]{
                \texttt{\tiny{[GM$\rangle$P2]}}
                \texttt{Kara: (5,5) nie jest puste, ale zawiera '│'.} \\
\\ 
\texttt{Łącznie zgromadziliście 16 z 16 kar.} \\
\texttt{Spróbuj ponownie!} \\
            }
        }
    }
    & \\ \\

    \theutterance \stepcounter{utterance}  
    & & & \multicolumn{4}{p{0.6\linewidth}}{
        \cellcolor[rgb]{0.9,0.9,0.9}{
            \makecell[{{p{\linewidth}}}]{
                \texttt{\tiny{[GM$\langle$P2]}}
                \texttt{RUCH: T, (4, 4)} \\
            }
        }
    }
     \\ \\

    \theutterance \stepcounter{utterance}  
    & & & \multicolumn{2}{p{0.3\linewidth}}{
        \cellcolor[rgb]{0.9,0.9,0.9}{
            \makecell[{{p{\linewidth}}}]{
                \texttt{\tiny{[GM$|$GM]}}
                \texttt{RUCH: T, (4, 4)} \\
            }
        }
    }
    & & \\ \\

    \theutterance \stepcounter{utterance}  
    & & & \multicolumn{2}{p{0.3\linewidth}}{
        \cellcolor[rgb]{0.9,0.9,0.9}{
            \makecell[{{p{\linewidth}}}]{
                \texttt{\tiny{[GM$|$GM]}}
                \texttt{Kara: (4,4) nie jest puste, ale zawiera '─'.} \\
\\ 
\texttt{Łącznie zgromadziliście 17 z 16 kar.} \\
\texttt{Spróbuj ponownie!} \\
            }
        }
    }
    & & \\ \\

    \theutterance \stepcounter{utterance}  
    & & \multicolumn{4}{p{0.6\linewidth}}{
        \cellcolor[rgb]{0.9,0.9,0.9}{
            \makecell[{{p{\linewidth}}}]{
                \texttt{\tiny{[GM$\rangle$P2]}}
                \texttt{Kara: (4,4) nie jest puste, ale zawiera '─'.} \\
\\ 
\texttt{Łącznie zgromadziliście 17 z 16 kar.} \\
\texttt{Spróbuj ponownie!} \\
            }
        }
    }
    & \\ \\

    \theutterance \stepcounter{utterance}  
    & & & \multicolumn{4}{p{0.6\linewidth}}{
        \cellcolor[rgb]{0.9,0.9,0.9}{
            \makecell[{{p{\linewidth}}}]{
                \texttt{\tiny{[GM$\langle$P2]}}
                \texttt{RUCH: T, (5, 4)} \\
            }
        }
    }
     \\ \\

    \theutterance \stepcounter{utterance}  
    & & & \multicolumn{2}{p{0.3\linewidth}}{
        \cellcolor[rgb]{0.9,0.9,0.9}{
            \makecell[{{p{\linewidth}}}]{
                \texttt{\tiny{[GM$|$GM]}}
                \texttt{RUCH: T, (5, 4)} \\
            }
        }
    }
    & & \\ \\

    \theutterance \stepcounter{utterance}  
    & & & \multicolumn{2}{p{0.3\linewidth}}{
        \cellcolor[rgb]{0.9,0.9,0.9}{
            \makecell[{{p{\linewidth}}}]{
                \texttt{\tiny{[GM$|$GM]}}
                \texttt{Kara: (5,4) nie jest puste, ale zawiera '┼'.} \\
\\ 
\texttt{Łącznie zgromadziliście 18 z 16 kar.} \\
\texttt{Spróbuj ponownie!} \\
            }
        }
    }
    & & \\ \\

    \theutterance \stepcounter{utterance}  
    & & \multicolumn{4}{p{0.6\linewidth}}{
        \cellcolor[rgb]{0.9,0.9,0.9}{
            \makecell[{{p{\linewidth}}}]{
                \texttt{\tiny{[GM$\rangle$P2]}}
                \texttt{Kara: (5,4) nie jest puste, ale zawiera '┼'.} \\
\\ 
\texttt{Łącznie zgromadziliście 18 z 16 kar.} \\
\texttt{Spróbuj ponownie!} \\
            }
        }
    }
    & \\ \\

    \theutterance \stepcounter{utterance}  
    & & & \multicolumn{4}{p{0.6\linewidth}}{
        \cellcolor[rgb]{0.9,0.9,0.9}{
            \makecell[{{p{\linewidth}}}]{
                \texttt{\tiny{[GM$\langle$P2]}}
                \texttt{RUCH: T, (5, 5)} \\
            }
        }
    }
     \\ \\

    \theutterance \stepcounter{utterance}  
    & & & \multicolumn{2}{p{0.3\linewidth}}{
        \cellcolor[rgb]{0.9,0.9,0.9}{
            \makecell[{{p{\linewidth}}}]{
                \texttt{\tiny{[GM$|$GM]}}
                \texttt{RUCH: T, (5, 5)} \\
            }
        }
    }
    & & \\ \\

    \theutterance \stepcounter{utterance}  
    & & & \multicolumn{2}{p{0.3\linewidth}}{
        \cellcolor[rgb]{0.9,0.9,0.9}{
            \makecell[{{p{\linewidth}}}]{
                \texttt{\tiny{[GM$|$GM]}}
                \texttt{Kara: (5,5) nie jest puste, ale zawiera '│'.} \\
\\ 
\texttt{Łącznie zgromadziliście 19 z 16 kar.} \\
\texttt{Spróbuj ponownie!} \\
            }
        }
    }
    & & \\ \\

    \theutterance \stepcounter{utterance}  
    & & \multicolumn{4}{p{0.6\linewidth}}{
        \cellcolor[rgb]{0.9,0.9,0.9}{
            \makecell[{{p{\linewidth}}}]{
                \texttt{\tiny{[GM$\rangle$P2]}}
                \texttt{Kara: (5,5) nie jest puste, ale zawiera '│'.} \\
\\ 
\texttt{Łącznie zgromadziliście 19 z 16 kar.} \\
\texttt{Spróbuj ponownie!} \\
            }
        }
    }
    & \\ \\

    \theutterance \stepcounter{utterance}  
    & & & \multicolumn{4}{p{0.6\linewidth}}{
        \cellcolor[rgb]{0.9,0.9,0.9}{
            \makecell[{{p{\linewidth}}}]{
                \texttt{\tiny{[GM$\langle$P2]}}
                \texttt{RUCH: T, (4, 4)} \\
            }
        }
    }
     \\ \\

    \theutterance \stepcounter{utterance}  
    & & & \multicolumn{2}{p{0.3\linewidth}}{
        \cellcolor[rgb]{0.9,0.9,0.9}{
            \makecell[{{p{\linewidth}}}]{
                \texttt{\tiny{[GM$|$GM]}}
                \texttt{RUCH: T, (4, 4)} \\
            }
        }
    }
    & & \\ \\

    \theutterance \stepcounter{utterance}  
    & & & \multicolumn{2}{p{0.3\linewidth}}{
        \cellcolor[rgb]{0.9,0.9,0.9}{
            \makecell[{{p{\linewidth}}}]{
                \texttt{\tiny{[GM$|$GM]}}
                \texttt{Kara: (4,4) nie jest puste, ale zawiera '─'.} \\
\\ 
\texttt{Łącznie zgromadziliście 20 z 16 kar.} \\
\texttt{Spróbuj ponownie!} \\
            }
        }
    }
    & & \\ \\

    \theutterance \stepcounter{utterance}  
    & & \multicolumn{4}{p{0.6\linewidth}}{
        \cellcolor[rgb]{0.9,0.9,0.9}{
            \makecell[{{p{\linewidth}}}]{
                \texttt{\tiny{[GM$\rangle$P2]}}
                \texttt{Kara: (4,4) nie jest puste, ale zawiera '─'.} \\
\\ 
\texttt{Łącznie zgromadziliście 20 z 16 kar.} \\
\texttt{Spróbuj ponownie!} \\
            }
        }
    }
    & \\ \\

    \theutterance \stepcounter{utterance}  
    & & & \multicolumn{4}{p{0.6\linewidth}}{
        \cellcolor[rgb]{0.9,0.9,0.9}{
            \makecell[{{p{\linewidth}}}]{
                \texttt{\tiny{[GM$\langle$P2]}}
                \texttt{RUCH: T, (4, 5)} \\
            }
        }
    }
     \\ \\

    \theutterance \stepcounter{utterance}  
    & & & \multicolumn{2}{p{0.3\linewidth}}{
        \cellcolor[rgb]{0.9,0.9,0.9}{
            \makecell[{{p{\linewidth}}}]{
                \texttt{\tiny{[GM$|$GM]}}
                \texttt{RUCH: T, (4, 5)} \\
            }
        }
    }
    & & \\ \\

    \theutterance \stepcounter{utterance}  
    & & & \multicolumn{2}{p{0.3\linewidth}}{
        \cellcolor[rgb]{0.9,0.9,0.9}{
            \makecell[{{p{\linewidth}}}]{
                \texttt{\tiny{[GM$|$GM]}}
                \texttt{T moved to (4, 5)} \\
            }
        }
    }
    & & \\ \\

    \theutterance \stepcounter{utterance}  
    & & \multicolumn{4}{p{0.6\linewidth}}{
        \cellcolor[rgb]{0.9,0.9,0.9}{
            \makecell[{{p{\linewidth}}}]{
                \texttt{\tiny{[P1$\langle$GM]}}
                \texttt{Pomyślnie przeniesiono 'S' do (3,3). Twoja zaktualizowana siatka wygląda następująco:} \\
\texttt{```} \\
\texttt{ 1234567} \\
\texttt{╔════╤══╗} \\
\texttt{║◌◌◌◌│◌◌║ 1} \\
\texttt{║P◌◌◌├──╢ 2} \\
\texttt{║◌◌SA│◌◌║ 3} \\
\texttt{║O┌──┼──╢ 4} \\
\texttt{║◌│T◌│◌D║ 5} \\
\texttt{╟─┴┬─┤◌M║ 6} \\
\texttt{║◌◌│◌│◌◌║ 7} \\
\texttt{╚══╧═╧══╝} \\
\texttt{```} \\
\texttt{Obecnie grasz rundę 13 z maksymalnie 28.} \\
\texttt{Łącznie zgromadziliście 20 z 16 kar.} \\
\texttt{Drugi gracz przesunął obiekt na swojej siatce.} \\
\\ 
\\ 
\texttt{Jakie jest twoje następne polecenie?} \\
            }
        }
    }
    & \\ \\

    \theutterance \stepcounter{utterance}  
    & \multicolumn{4}{p{0.6\linewidth}}{
        \cellcolor[rgb]{0.9,0.9,0.9}{
            \makecell[{{p{\linewidth}}}]{
                \texttt{\tiny{[P1$\rangle$GM]}}
                \texttt{RUCH: D, (5, 2)} \\
            }
        }
    }
    & & \\ \\

    \theutterance \stepcounter{utterance}  
    & & & \multicolumn{2}{p{0.3\linewidth}}{
        \cellcolor[rgb]{0.9,0.9,0.9}{
            \makecell[{{p{\linewidth}}}]{
                \texttt{\tiny{[GM$|$GM]}}
                \texttt{RUCH: D, (5, 2)} \\
            }
        }
    }
    & & \\ \\

    \theutterance \stepcounter{utterance}  
    & & & \multicolumn{2}{p{0.3\linewidth}}{
        \cellcolor[rgb]{0.9,0.9,0.9}{
            \makecell[{{p{\linewidth}}}]{
                \texttt{\tiny{[GM$|$GM]}}
                \texttt{Kara: (5,2) nie jest puste, ale zawiera '├'.} \\
\\ 
\texttt{Łącznie zgromadziliście 21 z 16 kar.} \\
\texttt{Spróbuj ponownie!} \\
            }
        }
    }
    & & \\ \\

    \theutterance \stepcounter{utterance}  
    & & \multicolumn{4}{p{0.6\linewidth}}{
        \cellcolor[rgb]{0.9,0.9,0.9}{
            \makecell[{{p{\linewidth}}}]{
                \texttt{\tiny{[P1$\langle$GM]}}
                \texttt{Kara: (5,2) nie jest puste, ale zawiera '├'.} \\
\\ 
\texttt{Łącznie zgromadziliście 21 z 16 kar.} \\
\texttt{Spróbuj ponownie!} \\
            }
        }
    }
    & \\ \\

    \theutterance \stepcounter{utterance}  
    & \multicolumn{4}{p{0.6\linewidth}}{
        \cellcolor[rgb]{0.9,0.9,0.9}{
            \makecell[{{p{\linewidth}}}]{
                \texttt{\tiny{[P1$\rangle$GM]}}
                \texttt{RUCH: D, (6, 2)} \\
            }
        }
    }
    & & \\ \\

    \theutterance \stepcounter{utterance}  
    & & & \multicolumn{2}{p{0.3\linewidth}}{
        \cellcolor[rgb]{0.9,0.9,0.9}{
            \makecell[{{p{\linewidth}}}]{
                \texttt{\tiny{[GM$|$GM]}}
                \texttt{RUCH: D, (6, 2)} \\
            }
        }
    }
    & & \\ \\

    \theutterance \stepcounter{utterance}  
    & & & \multicolumn{2}{p{0.3\linewidth}}{
        \cellcolor[rgb]{0.9,0.9,0.9}{
            \makecell[{{p{\linewidth}}}]{
                \texttt{\tiny{[GM$|$GM]}}
                \texttt{Kara: (6,2) nie jest puste, ale zawiera '─'.} \\
\\ 
\texttt{Łącznie zgromadziliście 22 z 16 kar.} \\
\texttt{Spróbuj ponownie!} \\
            }
        }
    }
    & & \\ \\

    \theutterance \stepcounter{utterance}  
    & & \multicolumn{4}{p{0.6\linewidth}}{
        \cellcolor[rgb]{0.9,0.9,0.9}{
            \makecell[{{p{\linewidth}}}]{
                \texttt{\tiny{[P1$\langle$GM]}}
                \texttt{Kara: (6,2) nie jest puste, ale zawiera '─'.} \\
\\ 
\texttt{Łącznie zgromadziliście 22 z 16 kar.} \\
\texttt{Spróbuj ponownie!} \\
            }
        }
    }
    & \\ \\

    \theutterance \stepcounter{utterance}  
    & \multicolumn{4}{p{0.6\linewidth}}{
        \cellcolor[rgb]{0.9,0.9,0.9}{
            \makecell[{{p{\linewidth}}}]{
                \texttt{\tiny{[P1$\rangle$GM]}}
                \texttt{RUCH: D, (4, 2)} \\
            }
        }
    }
    & & \\ \\

    \theutterance \stepcounter{utterance}  
    & & & \multicolumn{2}{p{0.3\linewidth}}{
        \cellcolor[rgb]{0.9,0.9,0.9}{
            \makecell[{{p{\linewidth}}}]{
                \texttt{\tiny{[GM$|$GM]}}
                \texttt{RUCH: D, (4, 2)} \\
            }
        }
    }
    & & \\ \\

    \theutterance \stepcounter{utterance}  
    & & & \multicolumn{2}{p{0.3\linewidth}}{
        \cellcolor[rgb]{0.9,0.9,0.9}{
            \makecell[{{p{\linewidth}}}]{
                \texttt{\tiny{[GM$|$GM]}}
                \texttt{D moved to (4, 2)} \\
            }
        }
    }
    & & \\ \\

    \theutterance \stepcounter{utterance}  
    & & \multicolumn{4}{p{0.6\linewidth}}{
        \cellcolor[rgb]{0.9,0.9,0.9}{
            \makecell[{{p{\linewidth}}}]{
                \texttt{\tiny{[GM$\rangle$P2]}}
                \texttt{Pomyślnie przeniesiono 'T' do (4,5). Twoja zaktualizowana siatka wygląda następująco:} \\
\texttt{```} \\
\texttt{ 1234567} \\
\texttt{╔════╤══╗} \\
\texttt{║◌P◌◌│◌◌║ 1} \\
\texttt{║◌◌AS├──╢ 2} \\
\texttt{║M◌◌◌│◌◌║ 3} \\
\texttt{║O┌──┼──╢ 4} \\
\texttt{║◌│◌T│D◌║ 5} \\
\texttt{╟─┴┬─┤◌◌║ 6} \\
\texttt{║◌◌│◌│◌◌║ 7} \\
\texttt{╚══╧═╧══╝} \\
\texttt{```} \\
\texttt{Obecnie grasz rundę 14 z maksymalnie 28.} \\
\texttt{Łącznie zgromadziliście 22 z 16 kar.} \\
\texttt{Drugi gracz przesunął obiekt na swojej siatce.} \\
\\ 
\\ 
\texttt{Jakie jest twoje następne polecenie?} \\
            }
        }
    }
    & \\ \\

    \theutterance \stepcounter{utterance}  
    & & & \multicolumn{4}{p{0.6\linewidth}}{
        \cellcolor[rgb]{0.9,0.9,0.9}{
            \makecell[{{p{\linewidth}}}]{
                \texttt{\tiny{[GM$\langle$P2]}}
                \texttt{RUCH: D, (5, 5)} \\
            }
        }
    }
     \\ \\

    \theutterance \stepcounter{utterance}  
    & & & \multicolumn{2}{p{0.3\linewidth}}{
        \cellcolor[rgb]{0.9,0.9,0.9}{
            \makecell[{{p{\linewidth}}}]{
                \texttt{\tiny{[GM$|$GM]}}
                \texttt{RUCH: D, (5, 5)} \\
            }
        }
    }
    & & \\ \\

    \theutterance \stepcounter{utterance}  
    & & & \multicolumn{2}{p{0.3\linewidth}}{
        \cellcolor[rgb]{0.9,0.9,0.9}{
            \makecell[{{p{\linewidth}}}]{
                \texttt{\tiny{[GM$|$GM]}}
                \texttt{Kara: (5,5) nie jest puste, ale zawiera '│'.} \\
\\ 
\texttt{Łącznie zgromadziliście 23 z 16 kar.} \\
\texttt{Spróbuj ponownie!} \\
            }
        }
    }
    & & \\ \\

    \theutterance \stepcounter{utterance}  
    & & \multicolumn{4}{p{0.6\linewidth}}{
        \cellcolor[rgb]{0.9,0.9,0.9}{
            \makecell[{{p{\linewidth}}}]{
                \texttt{\tiny{[GM$\rangle$P2]}}
                \texttt{Kara: (5,5) nie jest puste, ale zawiera '│'.} \\
\\ 
\texttt{Łącznie zgromadziliście 23 z 16 kar.} \\
\texttt{Spróbuj ponownie!} \\
            }
        }
    }
    & \\ \\

    \theutterance \stepcounter{utterance}  
    & & & \multicolumn{4}{p{0.6\linewidth}}{
        \cellcolor[rgb]{0.9,0.9,0.9}{
            \makecell[{{p{\linewidth}}}]{
                \texttt{\tiny{[GM$\langle$P2]}}
                \texttt{RUCH: D, (6, 5)} \\
            }
        }
    }
     \\ \\

    \theutterance \stepcounter{utterance}  
    & & & \multicolumn{2}{p{0.3\linewidth}}{
        \cellcolor[rgb]{0.9,0.9,0.9}{
            \makecell[{{p{\linewidth}}}]{
                \texttt{\tiny{[GM$|$GM]}}
                \texttt{RUCH: D, (6, 5)} \\
            }
        }
    }
    & & \\ \\

    \theutterance \stepcounter{utterance}  
    & & & \multicolumn{2}{p{0.3\linewidth}}{
        \cellcolor[rgb]{0.9,0.9,0.9}{
            \makecell[{{p{\linewidth}}}]{
                \texttt{\tiny{[GM$|$GM]}}
                \texttt{Kara: (6,5) nie jest puste, ale zawiera 'D'.} \\
\\ 
\texttt{Łącznie zgromadziliście 24 z 16 kar.} \\
\texttt{Spróbuj ponownie!} \\
            }
        }
    }
    & & \\ \\

    \theutterance \stepcounter{utterance}  
    & & \multicolumn{4}{p{0.6\linewidth}}{
        \cellcolor[rgb]{0.9,0.9,0.9}{
            \makecell[{{p{\linewidth}}}]{
                \texttt{\tiny{[GM$\rangle$P2]}}
                \texttt{Kara: (6,5) nie jest puste, ale zawiera 'D'.} \\
\\ 
\texttt{Łącznie zgromadziliście 24 z 16 kar.} \\
\texttt{Spróbuj ponownie!} \\
            }
        }
    }
    & \\ \\

    \theutterance \stepcounter{utterance}  
    & & & \multicolumn{4}{p{0.6\linewidth}}{
        \cellcolor[rgb]{0.9,0.9,0.9}{
            \makecell[{{p{\linewidth}}}]{
                \texttt{\tiny{[GM$\langle$P2]}}
                \texttt{RUCH: D, (6, 6)} \\
            }
        }
    }
     \\ \\

    \theutterance \stepcounter{utterance}  
    & & & \multicolumn{2}{p{0.3\linewidth}}{
        \cellcolor[rgb]{0.9,0.9,0.9}{
            \makecell[{{p{\linewidth}}}]{
                \texttt{\tiny{[GM$|$GM]}}
                \texttt{RUCH: D, (6, 6)} \\
            }
        }
    }
    & & \\ \\

    \theutterance \stepcounter{utterance}  
    & & & \multicolumn{2}{p{0.3\linewidth}}{
        \cellcolor[rgb]{0.9,0.9,0.9}{
            \makecell[{{p{\linewidth}}}]{
                \texttt{\tiny{[GM$|$GM]}}
                \texttt{D moved to (6, 6)} \\
            }
        }
    }
    & & \\ \\

    \theutterance \stepcounter{utterance}  
    & & \multicolumn{4}{p{0.6\linewidth}}{
        \cellcolor[rgb]{0.9,0.9,0.9}{
            \makecell[{{p{\linewidth}}}]{
                \texttt{\tiny{[P1$\langle$GM]}}
                \texttt{Pomyślnie przeniesiono 'D' do (4,2). Twoja zaktualizowana siatka wygląda następująco:} \\
\texttt{```} \\
\texttt{ 1234567} \\
\texttt{╔════╤══╗} \\
\texttt{║◌◌◌◌│◌◌║ 1} \\
\texttt{║P◌◌D├──╢ 2} \\
\texttt{║◌◌SA│◌◌║ 3} \\
\texttt{║O┌──┼──╢ 4} \\
\texttt{║◌│T◌│◌◌║ 5} \\
\texttt{╟─┴┬─┤◌M║ 6} \\
\texttt{║◌◌│◌│◌◌║ 7} \\
\texttt{╚══╧═╧══╝} \\
\texttt{```} \\
\texttt{Obecnie grasz rundę 14 z maksymalnie 28.} \\
\texttt{Łącznie zgromadziliście 24 z 16 kar.} \\
\texttt{Drugi gracz przesunął obiekt na swojej siatce.} \\
\\ 
\\ 
\texttt{Jakie jest twoje następne polecenie?} \\
            }
        }
    }
    & \\ \\

    \theutterance \stepcounter{utterance}  
    & \multicolumn{4}{p{0.6\linewidth}}{
        \cellcolor[rgb]{0.9,0.9,0.9}{
            \makecell[{{p{\linewidth}}}]{
                \texttt{\tiny{[P1$\rangle$GM]}}
                \texttt{RUCH: M, (5, 4)} \\
            }
        }
    }
    & & \\ \\

    \theutterance \stepcounter{utterance}  
    & & & \multicolumn{2}{p{0.3\linewidth}}{
        \cellcolor[rgb]{0.9,0.9,0.9}{
            \makecell[{{p{\linewidth}}}]{
                \texttt{\tiny{[GM$|$GM]}}
                \texttt{RUCH: M, (5, 4)} \\
            }
        }
    }
    & & \\ \\

    \theutterance \stepcounter{utterance}  
    & & & \multicolumn{2}{p{0.3\linewidth}}{
        \cellcolor[rgb]{0.9,0.9,0.9}{
            \makecell[{{p{\linewidth}}}]{
                \texttt{\tiny{[GM$|$GM]}}
                \texttt{Kara: (5,4) nie jest puste, ale zawiera '┼'.} \\
\\ 
\texttt{Łącznie zgromadziliście 25 z 16 kar.} \\
\texttt{Spróbuj ponownie!} \\
            }
        }
    }
    & & \\ \\

    \theutterance \stepcounter{utterance}  
    & & \multicolumn{4}{p{0.6\linewidth}}{
        \cellcolor[rgb]{0.9,0.9,0.9}{
            \makecell[{{p{\linewidth}}}]{
                \texttt{\tiny{[P1$\langle$GM]}}
                \texttt{Kara: (5,4) nie jest puste, ale zawiera '┼'.} \\
\\ 
\texttt{Łącznie zgromadziliście 25 z 16 kar.} \\
\texttt{Spróbuj ponownie!} \\
            }
        }
    }
    & \\ \\

    \theutterance \stepcounter{utterance}  
    & \multicolumn{4}{p{0.6\linewidth}}{
        \cellcolor[rgb]{0.9,0.9,0.9}{
            \makecell[{{p{\linewidth}}}]{
                \texttt{\tiny{[P1$\rangle$GM]}}
                \texttt{RUCH: M, (6, 4)} \\
            }
        }
    }
    & & \\ \\

    \theutterance \stepcounter{utterance}  
    & & & \multicolumn{2}{p{0.3\linewidth}}{
        \cellcolor[rgb]{0.9,0.9,0.9}{
            \makecell[{{p{\linewidth}}}]{
                \texttt{\tiny{[GM$|$GM]}}
                \texttt{RUCH: M, (6, 4)} \\
            }
        }
    }
    & & \\ \\

    \theutterance \stepcounter{utterance}  
    & & & \multicolumn{2}{p{0.3\linewidth}}{
        \cellcolor[rgb]{0.9,0.9,0.9}{
            \makecell[{{p{\linewidth}}}]{
                \texttt{\tiny{[GM$|$GM]}}
                \texttt{Kara: (6,4) nie jest puste, ale zawiera '─'.} \\
\\ 
\texttt{Łącznie zgromadziliście 26 z 16 kar.} \\
\texttt{Spróbuj ponownie!} \\
            }
        }
    }
    & & \\ \\

    \theutterance \stepcounter{utterance}  
    & & \multicolumn{4}{p{0.6\linewidth}}{
        \cellcolor[rgb]{0.9,0.9,0.9}{
            \makecell[{{p{\linewidth}}}]{
                \texttt{\tiny{[P1$\langle$GM]}}
                \texttt{Kara: (6,4) nie jest puste, ale zawiera '─'.} \\
\\ 
\texttt{Łącznie zgromadziliście 26 z 16 kar.} \\
\texttt{Spróbuj ponownie!} \\
            }
        }
    }
    & \\ \\

    \theutterance \stepcounter{utterance}  
    & \multicolumn{4}{p{0.6\linewidth}}{
        \cellcolor[rgb]{0.9,0.9,0.9}{
            \makecell[{{p{\linewidth}}}]{
                \texttt{\tiny{[P1$\rangle$GM]}}
                \texttt{RUCH: M, (5, 2)} \\
            }
        }
    }
    & & \\ \\

    \theutterance \stepcounter{utterance}  
    & & & \multicolumn{2}{p{0.3\linewidth}}{
        \cellcolor[rgb]{0.9,0.9,0.9}{
            \makecell[{{p{\linewidth}}}]{
                \texttt{\tiny{[GM$|$GM]}}
                \texttt{RUCH: M, (5, 2)} \\
            }
        }
    }
    & & \\ \\

    \theutterance \stepcounter{utterance}  
    & & & \multicolumn{2}{p{0.3\linewidth}}{
        \cellcolor[rgb]{0.9,0.9,0.9}{
            \makecell[{{p{\linewidth}}}]{
                \texttt{\tiny{[GM$|$GM]}}
                \texttt{Kara: (5,2) nie jest puste, ale zawiera '├'.} \\
\\ 
\texttt{Łącznie zgromadziliście 27 z 16 kar.} \\
\texttt{Spróbuj ponownie!} \\
            }
        }
    }
    & & \\ \\

    \theutterance \stepcounter{utterance}  
    & & \multicolumn{4}{p{0.6\linewidth}}{
        \cellcolor[rgb]{0.9,0.9,0.9}{
            \makecell[{{p{\linewidth}}}]{
                \texttt{\tiny{[P1$\langle$GM]}}
                \texttt{Kara: (5,2) nie jest puste, ale zawiera '├'.} \\
\\ 
\texttt{Łącznie zgromadziliście 27 z 16 kar.} \\
\texttt{Spróbuj ponownie!} \\
            }
        }
    }
    & \\ \\

    \theutterance \stepcounter{utterance}  
    & \multicolumn{4}{p{0.6\linewidth}}{
        \cellcolor[rgb]{0.9,0.9,0.9}{
            \makecell[{{p{\linewidth}}}]{
                \texttt{\tiny{[P1$\rangle$GM]}}
                \texttt{RUCH: M, (4, 3)} \\
            }
        }
    }
    & & \\ \\

    \theutterance \stepcounter{utterance}  
    & & & \multicolumn{2}{p{0.3\linewidth}}{
        \cellcolor[rgb]{0.9,0.9,0.9}{
            \makecell[{{p{\linewidth}}}]{
                \texttt{\tiny{[GM$|$GM]}}
                \texttt{RUCH: M, (4, 3)} \\
            }
        }
    }
    & & \\ \\

    \theutterance \stepcounter{utterance}  
    & & & \multicolumn{2}{p{0.3\linewidth}}{
        \cellcolor[rgb]{0.9,0.9,0.9}{
            \makecell[{{p{\linewidth}}}]{
                \texttt{\tiny{[GM$|$GM]}}
                \texttt{Kara: (4,3) nie jest puste, ale zawiera 'A'.} \\
\\ 
\texttt{Łącznie zgromadziliście 28 z 16 kar.} \\
\texttt{Spróbuj ponownie!} \\
            }
        }
    }
    & & \\ \\

    \theutterance \stepcounter{utterance}  
    & & \multicolumn{4}{p{0.6\linewidth}}{
        \cellcolor[rgb]{0.9,0.9,0.9}{
            \makecell[{{p{\linewidth}}}]{
                \texttt{\tiny{[P1$\langle$GM]}}
                \texttt{Kara: (4,3) nie jest puste, ale zawiera 'A'.} \\
\\ 
\texttt{Łącznie zgromadziliście 28 z 16 kar.} \\
\texttt{Spróbuj ponownie!} \\
            }
        }
    }
    & \\ \\

    \theutterance \stepcounter{utterance}  
    & \multicolumn{4}{p{0.6\linewidth}}{
        \cellcolor[rgb]{0.9,0.9,0.9}{
            \makecell[{{p{\linewidth}}}]{
                \texttt{\tiny{[P1$\rangle$GM]}}
                \texttt{RUCH: A, (3, 3)} \\
            }
        }
    }
    & & \\ \\

    \theutterance \stepcounter{utterance}  
    & & & \multicolumn{2}{p{0.3\linewidth}}{
        \cellcolor[rgb]{0.9,0.9,0.9}{
            \makecell[{{p{\linewidth}}}]{
                \texttt{\tiny{[GM$|$GM]}}
                \texttt{RUCH: A, (3, 3)} \\
            }
        }
    }
    & & \\ \\

    \theutterance \stepcounter{utterance}  
    & & & \multicolumn{2}{p{0.3\linewidth}}{
        \cellcolor[rgb]{0.9,0.9,0.9}{
            \makecell[{{p{\linewidth}}}]{
                \texttt{\tiny{[GM$|$GM]}}
                \texttt{Kara: (3,3) nie jest puste, ale zawiera 'S'.} \\
\\ 
\texttt{Łącznie zgromadziliście 29 z 16 kar.} \\
\texttt{Spróbuj ponownie!} \\
            }
        }
    }
    & & \\ \\

    \theutterance \stepcounter{utterance}  
    & & \multicolumn{4}{p{0.6\linewidth}}{
        \cellcolor[rgb]{0.9,0.9,0.9}{
            \makecell[{{p{\linewidth}}}]{
                \texttt{\tiny{[P1$\langle$GM]}}
                \texttt{Kara: (3,3) nie jest puste, ale zawiera 'S'.} \\
\\ 
\texttt{Łącznie zgromadziliście 29 z 16 kar.} \\
\texttt{Spróbuj ponownie!} \\
            }
        }
    }
    & \\ \\

    \theutterance \stepcounter{utterance}  
    & \multicolumn{4}{p{0.6\linewidth}}{
        \cellcolor[rgb]{0.9,0.9,0.9}{
            \makecell[{{p{\linewidth}}}]{
                \texttt{\tiny{[P1$\rangle$GM]}}
                \texttt{RUCH: A, (4, 1)} \\
            }
        }
    }
    & & \\ \\

    \theutterance \stepcounter{utterance}  
    & & & \multicolumn{2}{p{0.3\linewidth}}{
        \cellcolor[rgb]{0.9,0.9,0.9}{
            \makecell[{{p{\linewidth}}}]{
                \texttt{\tiny{[GM$|$GM]}}
                \texttt{RUCH: A, (4, 1)} \\
            }
        }
    }
    & & \\ \\

    \theutterance \stepcounter{utterance}  
    & & & \multicolumn{2}{p{0.3\linewidth}}{
        \cellcolor[rgb]{0.9,0.9,0.9}{
            \makecell[{{p{\linewidth}}}]{
                \texttt{\tiny{[GM$|$GM]}}
                \texttt{A moved to (4, 1)} \\
            }
        }
    }
    & & \\ \\

    \theutterance \stepcounter{utterance}  
    & & \multicolumn{4}{p{0.6\linewidth}}{
        \cellcolor[rgb]{0.9,0.9,0.9}{
            \makecell[{{p{\linewidth}}}]{
                \texttt{\tiny{[GM$\rangle$P2]}}
                \texttt{Pomyślnie przeniesiono 'D' do (6,6). Twoja zaktualizowana siatka wygląda następująco:} \\
\texttt{```} \\
\texttt{ 1234567} \\
\texttt{╔════╤══╗} \\
\texttt{║◌P◌◌│◌◌║ 1} \\
\texttt{║◌◌AS├──╢ 2} \\
\texttt{║M◌◌◌│◌◌║ 3} \\
\texttt{║O┌──┼──╢ 4} \\
\texttt{║◌│◌T│◌◌║ 5} \\
\texttt{╟─┴┬─┤D◌║ 6} \\
\texttt{║◌◌│◌│◌◌║ 7} \\
\texttt{╚══╧═╧══╝} \\
\texttt{```} \\
\texttt{Obecnie grasz rundę 15 z maksymalnie 28.} \\
\texttt{Łącznie zgromadziliście 29 z 16 kar.} \\
\texttt{Drugi gracz przesunął obiekt na swojej siatce.} \\
\\ 
\\ 
\texttt{Jakie jest twoje następne polecenie?} \\
            }
        }
    }
    & \\ \\

    \theutterance \stepcounter{utterance}  
    & & & \multicolumn{4}{p{0.6\linewidth}}{
        \cellcolor[rgb]{0.9,0.9,0.9}{
            \makecell[{{p{\linewidth}}}]{
                \texttt{\tiny{[GM$\langle$P2]}}
                \texttt{RUCH: O, (3, 4)} \\
            }
        }
    }
     \\ \\

    \theutterance \stepcounter{utterance}  
    & & & \multicolumn{2}{p{0.3\linewidth}}{
        \cellcolor[rgb]{0.9,0.9,0.9}{
            \makecell[{{p{\linewidth}}}]{
                \texttt{\tiny{[GM$|$GM]}}
                \texttt{RUCH: O, (3, 4)} \\
            }
        }
    }
    & & \\ \\

    \theutterance \stepcounter{utterance}  
    & & & \multicolumn{2}{p{0.3\linewidth}}{
        \cellcolor[rgb]{0.9,0.9,0.9}{
            \makecell[{{p{\linewidth}}}]{
                \texttt{\tiny{[GM$|$GM]}}
                \texttt{Kara: (3,4) nie jest puste, ale zawiera '─'.} \\
\\ 
\texttt{Łącznie zgromadziliście 30 z 16 kar.} \\
\texttt{Spróbuj ponownie!} \\
            }
        }
    }
    & & \\ \\

    \theutterance \stepcounter{utterance}  
    & & \multicolumn{4}{p{0.6\linewidth}}{
        \cellcolor[rgb]{0.9,0.9,0.9}{
            \makecell[{{p{\linewidth}}}]{
                \texttt{\tiny{[GM$\rangle$P2]}}
                \texttt{Kara: (3,4) nie jest puste, ale zawiera '─'.} \\
\\ 
\texttt{Łącznie zgromadziliście 30 z 16 kar.} \\
\texttt{Spróbuj ponownie!} \\
            }
        }
    }
    & \\ \\

    \theutterance \stepcounter{utterance}  
    & & & \multicolumn{4}{p{0.6\linewidth}}{
        \cellcolor[rgb]{0.9,0.9,0.9}{
            \makecell[{{p{\linewidth}}}]{
                \texttt{\tiny{[GM$\langle$P2]}}
                \texttt{RUCH: O, (4, 4)} \\
            }
        }
    }
     \\ \\

    \theutterance \stepcounter{utterance}  
    & & & \multicolumn{2}{p{0.3\linewidth}}{
        \cellcolor[rgb]{0.9,0.9,0.9}{
            \makecell[{{p{\linewidth}}}]{
                \texttt{\tiny{[GM$|$GM]}}
                \texttt{RUCH: O, (4, 4)} \\
            }
        }
    }
    & & \\ \\

    \theutterance \stepcounter{utterance}  
    & & & \multicolumn{2}{p{0.3\linewidth}}{
        \cellcolor[rgb]{0.9,0.9,0.9}{
            \makecell[{{p{\linewidth}}}]{
                \texttt{\tiny{[GM$|$GM]}}
                \texttt{Kara: (4,4) nie jest puste, ale zawiera '─'.} \\
\\ 
\texttt{Łącznie zgromadziliście 31 z 16 kar.} \\
\texttt{Spróbuj ponownie!} \\
            }
        }
    }
    & & \\ \\

    \theutterance \stepcounter{utterance}  
    & & \multicolumn{4}{p{0.6\linewidth}}{
        \cellcolor[rgb]{0.9,0.9,0.9}{
            \makecell[{{p{\linewidth}}}]{
                \texttt{\tiny{[GM$\rangle$P2]}}
                \texttt{Kara: (4,4) nie jest puste, ale zawiera '─'.} \\
\\ 
\texttt{Łącznie zgromadziliście 31 z 16 kar.} \\
\texttt{Spróbuj ponownie!} \\
            }
        }
    }
    & \\ \\

    \theutterance \stepcounter{utterance}  
    & & & \multicolumn{4}{p{0.6\linewidth}}{
        \cellcolor[rgb]{0.9,0.9,0.9}{
            \makecell[{{p{\linewidth}}}]{
                \texttt{\tiny{[GM$\langle$P2]}}
                \texttt{RUCH: O, (4, 5)} \\
            }
        }
    }
     \\ \\

    \theutterance \stepcounter{utterance}  
    & & & \multicolumn{2}{p{0.3\linewidth}}{
        \cellcolor[rgb]{0.9,0.9,0.9}{
            \makecell[{{p{\linewidth}}}]{
                \texttt{\tiny{[GM$|$GM]}}
                \texttt{RUCH: O, (4, 5)} \\
            }
        }
    }
    & & \\ \\

    \theutterance \stepcounter{utterance}  
    & & & \multicolumn{2}{p{0.3\linewidth}}{
        \cellcolor[rgb]{0.9,0.9,0.9}{
            \makecell[{{p{\linewidth}}}]{
                \texttt{\tiny{[GM$|$GM]}}
                \texttt{Kara: (4,5) nie jest puste, ale zawiera 'T'.} \\
\\ 
\texttt{Łącznie zgromadziliście 32 z 16 kar.} \\
\texttt{Spróbuj ponownie!} \\
            }
        }
    }
    & & \\ \\

    \theutterance \stepcounter{utterance}  
    & & \multicolumn{4}{p{0.6\linewidth}}{
        \cellcolor[rgb]{0.9,0.9,0.9}{
            \makecell[{{p{\linewidth}}}]{
                \texttt{\tiny{[GM$\rangle$P2]}}
                \texttt{Kara: (4,5) nie jest puste, ale zawiera 'T'.} \\
\\ 
\texttt{Łącznie zgromadziliście 32 z 16 kar.} \\
\texttt{Spróbuj ponownie!} \\
            }
        }
    }
    & \\ \\

    \theutterance \stepcounter{utterance}  
    & & & \multicolumn{4}{p{0.6\linewidth}}{
        \cellcolor[rgb]{0.9,0.9,0.9}{
            \makecell[{{p{\linewidth}}}]{
                \texttt{\tiny{[GM$\langle$P2]}}
                \texttt{RUCH: O, (3, 4)} \\
            }
        }
    }
     \\ \\

    \theutterance \stepcounter{utterance}  
    & & & \multicolumn{2}{p{0.3\linewidth}}{
        \cellcolor[rgb]{0.9,0.9,0.9}{
            \makecell[{{p{\linewidth}}}]{
                \texttt{\tiny{[GM$|$GM]}}
                \texttt{RUCH: O, (3, 4)} \\
            }
        }
    }
    & & \\ \\

    \theutterance \stepcounter{utterance}  
    & & & \multicolumn{2}{p{0.3\linewidth}}{
        \cellcolor[rgb]{0.9,0.9,0.9}{
            \makecell[{{p{\linewidth}}}]{
                \texttt{\tiny{[GM$|$GM]}}
                \texttt{Kara: (3,4) nie jest puste, ale zawiera '─'.} \\
\\ 
\texttt{Łącznie zgromadziliście 33 z 16 kar.} \\
\texttt{Spróbuj ponownie!} \\
            }
        }
    }
    & & \\ \\

    \theutterance \stepcounter{utterance}  
    & & \multicolumn{4}{p{0.6\linewidth}}{
        \cellcolor[rgb]{0.9,0.9,0.9}{
            \makecell[{{p{\linewidth}}}]{
                \texttt{\tiny{[GM$\rangle$P2]}}
                \texttt{Kara: (3,4) nie jest puste, ale zawiera '─'.} \\
\\ 
\texttt{Łącznie zgromadziliście 33 z 16 kar.} \\
\texttt{Spróbuj ponownie!} \\
            }
        }
    }
    & \\ \\

    \theutterance \stepcounter{utterance}  
    & & & \multicolumn{4}{p{0.6\linewidth}}{
        \cellcolor[rgb]{0.9,0.9,0.9}{
            \makecell[{{p{\linewidth}}}]{
                \texttt{\tiny{[GM$\langle$P2]}}
                \texttt{RUCH: O, (4, 3)} \\
            }
        }
    }
     \\ \\

    \theutterance \stepcounter{utterance}  
    & & & \multicolumn{2}{p{0.3\linewidth}}{
        \cellcolor[rgb]{0.9,0.9,0.9}{
            \makecell[{{p{\linewidth}}}]{
                \texttt{\tiny{[GM$|$GM]}}
                \texttt{RUCH: O, (4, 3)} \\
            }
        }
    }
    & & \\ \\

    \theutterance \stepcounter{utterance}  
    & & & \multicolumn{2}{p{0.3\linewidth}}{
        \cellcolor[rgb]{0.9,0.9,0.9}{
            \makecell[{{p{\linewidth}}}]{
                \texttt{\tiny{[GM$|$GM]}}
                \texttt{O moved to (4, 3)} \\
            }
        }
    }
    & & \\ \\

    \theutterance \stepcounter{utterance}  
    & & \multicolumn{4}{p{0.6\linewidth}}{
        \cellcolor[rgb]{0.9,0.9,0.9}{
            \makecell[{{p{\linewidth}}}]{
                \texttt{\tiny{[P1$\langle$GM]}}
                \texttt{Pomyślnie przeniesiono 'A' do (4,1). Twoja zaktualizowana siatka wygląda następująco:} \\
\texttt{```} \\
\texttt{ 1234567} \\
\texttt{╔════╤══╗} \\
\texttt{║◌◌◌A│◌◌║ 1} \\
\texttt{║P◌◌D├──╢ 2} \\
\texttt{║◌◌S◌│◌◌║ 3} \\
\texttt{║O┌──┼──╢ 4} \\
\texttt{║◌│T◌│◌◌║ 5} \\
\texttt{╟─┴┬─┤◌M║ 6} \\
\texttt{║◌◌│◌│◌◌║ 7} \\
\texttt{╚══╧═╧══╝} \\
\texttt{```} \\
\texttt{Obecnie grasz rundę 15 z maksymalnie 28.} \\
\texttt{Łącznie zgromadziliście 33 z 16 kar.} \\
\texttt{Drugi gracz przesunął obiekt na swojej siatce.} \\
\\ 
\\ 
\texttt{Jakie jest twoje następne polecenie?} \\
            }
        }
    }
    & \\ \\

    \theutterance \stepcounter{utterance}  
    & \multicolumn{4}{p{0.6\linewidth}}{
        \cellcolor[rgb]{0.9,0.9,0.9}{
            \makecell[{{p{\linewidth}}}]{
                \texttt{\tiny{[P1$\rangle$GM]}}
                \texttt{POWIEDZ: Czy możemy teraz spróbować zakończyć grę, ponieważ wydaje się, że zbliżamy się do wyrównania wszystkich obiektów na obu planszach?} \\
            }
        }
    }
    & & \\ \\

    \theutterance \stepcounter{utterance}  
    & & & \multicolumn{2}{p{0.3\linewidth}}{
        \cellcolor[rgb]{0.9,0.9,0.9}{
            \makecell[{{p{\linewidth}}}]{
                \texttt{\tiny{[GM$|$GM]}}
                \texttt{POWIEDZ: Czy możemy teraz spróbować zakończyć grę, ponieważ wydaje się, że zbliżamy się do wyrównania wszystkich obiektów na obu planszach?} \\
            }
        }
    }
    & & \\ \\

    \theutterance \stepcounter{utterance}  
    & & \multicolumn{4}{p{0.6\linewidth}}{
        \cellcolor[rgb]{0.9,0.9,0.9}{
            \makecell[{{p{\linewidth}}}]{
                \texttt{\tiny{[GM$\rangle$P2]}}
                \texttt{Pomyślnie przeniesiono 'O' do (4,3). Twoja zaktualizowana siatka wygląda następująco:} \\
\texttt{```} \\
\texttt{ 1234567} \\
\texttt{╔════╤══╗} \\
\texttt{║◌P◌◌│◌◌║ 1} \\
\texttt{║◌◌AS├──╢ 2} \\
\texttt{║M◌◌O│◌◌║ 3} \\
\texttt{║◌┌──┼──╢ 4} \\
\texttt{║◌│◌T│◌◌║ 5} \\
\texttt{╟─┴┬─┤D◌║ 6} \\
\texttt{║◌◌│◌│◌◌║ 7} \\
\texttt{╚══╧═╧══╝} \\
\texttt{```} \\
\texttt{Obecnie grasz rundę 16 z maksymalnie 28.} \\
\texttt{Łącznie zgromadziliście 33 z 16 kar.} \\
\texttt{Drugi gracz wysłał tę wiadomość:} \\
\\ 
\\ 
\texttt{"Czy możemy teraz spróbować zakończyć grę, ponieważ wydaje się, że zbliżamy się do wyrównania wszystkich obiektów na obu planszach?"} \\
\\ 
\\ 
\texttt{Jakie jest twoje następne polecenie?} \\
            }
        }
    }
    & \\ \\

    \theutterance \stepcounter{utterance}  
    & & & \multicolumn{4}{p{0.6\linewidth}}{
        \cellcolor[rgb]{0.9,0.9,0.9}{
            \makecell[{{p{\linewidth}}}]{
                \texttt{\tiny{[GM$\langle$P2]}}
                \texttt{POWIEDZ: zakończono?} \\
            }
        }
    }
     \\ \\

    \theutterance \stepcounter{utterance}  
    & & & \multicolumn{2}{p{0.3\linewidth}}{
        \cellcolor[rgb]{0.9,0.9,0.9}{
            \makecell[{{p{\linewidth}}}]{
                \texttt{\tiny{[GM$|$GM]}}
                \texttt{POWIEDZ: zakończono?} \\
            }
        }
    }
    & & \\ \\

    \theutterance \stepcounter{utterance}  
    & & \multicolumn{4}{p{0.6\linewidth}}{
        \cellcolor[rgb]{0.9,0.9,0.9}{
            \makecell[{{p{\linewidth}}}]{
                \texttt{\tiny{[P1$\langle$GM]}}
                \texttt{Twoja wiadomość została przekazana drugiemu graczowi.} \\
\\ 
\\ 
\texttt{Obecnie grasz rundę 16 z maksymalnie 28.} \\
\texttt{Łącznie zgromadziliście 33 z 16 kar.} \\
\texttt{Drugi gracz wysłał tę wiadomość:} \\
\\ 
\\ 
\texttt{"zakończono?"} \\
\\ 
\\ 
\texttt{Jakie jest twoje następne polecenie?} \\
            }
        }
    }
    & \\ \\

    \theutterance \stepcounter{utterance}  
    & \multicolumn{4}{p{0.6\linewidth}}{
        \cellcolor[rgb]{0.9,0.9,0.9}{
            \makecell[{{p{\linewidth}}}]{
                \texttt{\tiny{[P1$\rangle$GM]}}
                \texttt{POWIEDZ: zakończono!} \\
            }
        }
    }
    & & \\ \\

    \theutterance \stepcounter{utterance}  
    & & & \multicolumn{2}{p{0.3\linewidth}}{
        \cellcolor[rgb]{0.9,0.9,0.9}{
            \makecell[{{p{\linewidth}}}]{
                \texttt{\tiny{[GM$|$GM]}}
                \texttt{POWIEDZ: zakończono!} \\
            }
        }
    }
    & & \\ \\

    \theutterance \stepcounter{utterance}  
    & & & \multicolumn{2}{p{0.3\linewidth}}{
        \cellcolor[rgb]{0.9,0.9,0.9}{
            \makecell[{{p{\linewidth}}}]{
                \texttt{\tiny{[GM$|$GM]}}
                \texttt{true} \\
            }
        }
    }
    & & \\ \\

    \theutterance \stepcounter{utterance}  
    & & & \multicolumn{2}{p{0.3\linewidth}}{
        \cellcolor[rgb]{0.9,0.9,0.9}{
            \makecell[{{p{\linewidth}}}]{
                \texttt{\tiny{[GM$|$GM]}}
                \texttt{Initial states:} \\
\texttt{```} \\
\texttt{Player 1:} \\
\texttt{ 1234567} \\
\texttt{╔════╤══╗} \\
\texttt{║◌◌◌O│PM║ 1} \\
\texttt{║◌◌T◌├──╢ 2} \\
\texttt{║◌◌◌◌│◌◌║ 3} \\
\texttt{║S┌──┼──╢ 4} \\
\texttt{║◌│◌A│◌◌║ 5} \\
\texttt{╟─┴┬─┤◌◌║ 6} \\
\texttt{║◌◌│◌│D◌║ 7} \\
\texttt{╚══╧═╧══╝} \\
\\ 
\\ 
\texttt{Player 2:} \\
\texttt{ 1234567} \\
\texttt{╔════╤══╗} \\
\texttt{║◌◌◌◌│O◌║ 1} \\
\texttt{║P◌M◌├──╢ 2} \\
\texttt{║◌S◌◌│T◌║ 3} \\
\texttt{║◌┌──┼──╢ 4} \\
\texttt{║A│◌◌│◌◌║ 5} \\
\texttt{╟─┴┬─┤◌D║ 6} \\
\texttt{║◌◌│◌│◌◌║ 7} \\
\texttt{╚══╧═╧══╝} \\
\\ 
\texttt{```} \\
            }
        }
    }
    & & \\ \\

    \theutterance \stepcounter{utterance}  
    & & & \multicolumn{2}{p{0.3\linewidth}}{
        \cellcolor[rgb]{0.9,0.9,0.9}{
            \makecell[{{p{\linewidth}}}]{
                \texttt{\tiny{[GM$|$GM]}}
                \texttt{End states:} \\
\texttt{```} \\
\texttt{Player 1:} \\
\texttt{ 1234567} \\
\texttt{╔════╤══╗} \\
\texttt{║◌◌◌A│◌◌║ 1} \\
\texttt{║P◌◌D├──╢ 2} \\
\texttt{║◌◌S◌│◌◌║ 3} \\
\texttt{║O┌──┼──╢ 4} \\
\texttt{║◌│T◌│◌◌║ 5} \\
\texttt{╟─┴┬─┤◌M║ 6} \\
\texttt{║◌◌│◌│◌◌║ 7} \\
\texttt{╚══╧═╧══╝} \\
\\ 
\\ 
\texttt{Player 2:} \\
\texttt{ 1234567} \\
\texttt{╔════╤══╗} \\
\texttt{║◌P◌◌│◌◌║ 1} \\
\texttt{║◌◌AS├──╢ 2} \\
\texttt{║M◌◌O│◌◌║ 3} \\
\texttt{║◌┌──┼──╢ 4} \\
\texttt{║◌│◌T│◌◌║ 5} \\
\texttt{╟─┴┬─┤D◌║ 6} \\
\texttt{║◌◌│◌│◌◌║ 7} \\
\texttt{╚══╧═╧══╝} \\
\\ 
\texttt{```} \\
            }
        }
    }
    & & \\ \\

    \theutterance \stepcounter{utterance}  
    & & & \multicolumn{2}{p{0.3\linewidth}}{
        \cellcolor[rgb]{0.9,0.9,0.9}{
            \makecell[{{p{\linewidth}}}]{
                \texttt{\tiny{[GM$|$GM]}}
                \texttt{* success: True} \\
\texttt{* lose: False} \\
\texttt{* aborted: False} \\
\texttt{{-}{-}{-}{-}{-}{-}{-}} \\
\texttt{* Shifts: 18.00} \\
\texttt{* Max Shifts: 12.00} \\
\texttt{* Min Shifts: 6.00} \\
\texttt{* End Distance Sum: 19.59} \\
\texttt{* Init Distance Sum: 20.21} \\
\texttt{* Expected Distance Sum: 29.33} \\
\texttt{* Penalties: 33.00} \\
\texttt{* Max Penalties: 16.00} \\
\texttt{* Rounds: 16.00} \\
\texttt{* Max Rounds: 28.00} \\
\texttt{* Object Count: 7.00} \\
            }
        }
    }
    & & \\ \\

\end{supertabular}
}

\end{document}
